\documentclass[12pt,letterpaper]{article}

%% ─── Packages ────────────────────────────────────────────────────────────────
\usepackage[margin=1in]{geometry}
\usepackage{amsmath,amssymb,amsthm}
\usepackage{mathtools}
\usepackage{hyperref}
\usepackage{xcolor}
\usepackage{listings}
\usepackage{booktabs}
\usepackage{array}
\usepackage{longtable}
\usepackage{graphicx}
\usepackage{fancyhdr}
\usepackage{titlesec}
\usepackage{enumitem}
\usepackage{microtype}
\usepackage{abstract}
\usepackage{caption}
\usepackage{subcaption}
\usepackage{tcolorbox}
\usepackage{mdframed}
\usepackage{algorithm}
\usepackage{algpseudocode}
\usepackage{cleveref}
\usepackage{soul}
\usepackage{setspace}
\usepackage{natbib}
\usepackage{float}

%% ─── Hyperref config ─────────────────────────────────────────────────────────
\hypersetup{
  colorlinks=true,
  linkcolor=blue!60!black,
  urlcolor=blue!70!black,
  citecolor=green!50!black,
  pdftitle={Gate U: Multiplicity-Native Hardware Acceleration --- Defensive Publication},
  pdfauthor={Multiplicity Foundation / PhaseMirror-HQ},
  pdfsubject={Prime-Indexed Tensor Acceleration, Governance-Aware ASIC Design},
  pdfkeywords={Multiplicity Theory, PIRTM, prime-indexed, FPGA, ASIC, MLIR, contractivity, governance}
}

%% ─── Listings config ─────────────────────────────────────────────────────────
\definecolor{codegreen}{rgb}{0,0.5,0}
\definecolor{codegray}{rgb}{0.5,0.5,0.5}
\definecolor{codepurple}{rgb}{0.58,0,0.82}
\definecolor{backcolour}{rgb}{0.97,0.97,0.97}
\definecolor{primepink}{rgb}{0.75,0.0,0.25}

\lstdefinestyle{pythonstyle}{
  backgroundcolor=\color{backcolour},
  commentstyle=\color{codegreen}\itshape,
  keywordstyle=\color{blue}\bfseries,
  numberstyle=\tiny\color{codegray},
  stringstyle=\color{codepurple},
  basicstyle=\ttfamily\small,
  breakatwhitespace=false,
  breaklines=true,
  captionpos=b,
  keepspaces=true,
  numbers=left,
  numbersep=5pt,
  showspaces=false,
  showstringspaces=false,
  showtabs=false,
  tabsize=2,
  frame=single,
  rulecolor=\color{gray!40},
  language=Python
}

\lstdefinestyle{cppstyle}{
  backgroundcolor=\color{backcolour},
  commentstyle=\color{codegreen}\itshape,
  keywordstyle=\color{blue!80!black}\bfseries,
  numberstyle=\tiny\color{codegray},
  stringstyle=\color{codepurple},
  basicstyle=\ttfamily\small,
  breakatwhitespace=false,
  breaklines=true,
  captionpos=b,
  keepspaces=true,
  numbers=left,
  numbersep=5pt,
  showspaces=false,
  showstringspaces=false,
  showtabs=false,
  tabsize=2,
  frame=single,
  rulecolor=\color{gray!40},
  language=C++
}

\lstdefinestyle{verilogstyle}{
  backgroundcolor=\color{backcolour},
  commentstyle=\color{codegreen}\itshape,
  keywordstyle=\color{primepink}\bfseries,
  numberstyle=\tiny\color{codegray},
  stringstyle=\color{codepurple},
  basicstyle=\ttfamily\small,
  breakatwhitespace=false,
  breaklines=true,
  captionpos=b,
  keepspaces=true,
  numbers=left,
  numbersep=5pt,
  showspaces=false,
  showstringspaces=false,
  showtabs=false,
  tabsize=2,
  frame=single,
  rulecolor=\color{gray!40},
  language=Verilog
}

\lstdefinestyle{tablegen}{
  backgroundcolor=\color{backcolour},
  commentstyle=\color{codegreen}\itshape,
  keywordstyle=\color{blue!80!black}\bfseries,
  numberstyle=\tiny\color{codegray},
  stringstyle=\color{codepurple},
  basicstyle=\ttfamily\small,
  breakatwhitespace=false,
  breaklines=true,
  captionpos=b,
  keepspaces=true,
  numbers=left,
  numbersep=5pt,
  showspaces=false,
  showstringspaces=false,
  showtabs=false,
  tabsize=2,
  frame=single,
  rulecolor=\color{gray!40},
  language={[Objective]C}
}

%% ─── Theorem environments ────────────────────────────────────────────────────
\theoremstyle{plain}
\newtheorem{theorem}{Theorem}[section]
\newtheorem{lemma}[theorem]{Lemma}
\newtheorem{proposition}[theorem]{Proposition}
\newtheorem{corollary}[theorem]{Corollary}

\theoremstyle{definition}
\newtheorem{definition}[theorem]{Definition}
\newtheorem{example}[theorem]{Example}
\newtheorem{invariant}[theorem]{Invariant}
\newtheorem{contract}[theorem]{Contract}

\theoremstyle{remark}
\newtheorem{remark}[theorem]{Remark}
\newtheorem{note}[theorem]{Note}

%% ─── tcolorbox environments ─────────────────────────────────────────────────
\tcbuselibrary{skins,breakable}

\newtcolorbox{defensebox}{
  colback=blue!4!white,
  colframe=blue!60!black,
  title={\textbf{Defensive Publication Notice}},
  fonttitle=\bfseries,
  breakable
}

\newtcolorbox{warningbox}{
  colback=red!5!white,
  colframe=red!60!black,
  title={\textbf{Open Tension / Known Limitation}},
  fonttitle=\bfseries,
  breakable
}

\newtcolorbox{invariantbox}{
  colback=green!4!white,
  colframe=green!50!black,
  title={\textbf{Normative Invariant}},
  fonttitle=\bfseries,
  breakable
}

%% ─── Headers and footers ─────────────────────────────────────────────────────
\pagestyle{fancy}
\fancyhf{}
\fancyhead[L]{\small\textit{Gate U: Multiplicity-Native Hardware Acceleration}}
\fancyhead[R]{\small\textit{Defensive Publication --- PhaseMirror-HQ}}
\fancyfoot[C]{\thepage}
\renewcommand{\headrulewidth}{0.4pt}

%% ─── Title config ─────────────────────────────────────────────────────────────
\title{
  \vspace{-1cm}
  {\Large \textbf{Defensive Publication}} \\[0.4em]
  {\huge \textbf{Gate U: Multiplicity-Native Hardware Acceleration}} \\[0.4em]
  {\Large Prime-Indexed Tensor Engines, Governance-Aware ASIC Design,\\
   and the \texttt{pirtm.step} Hardware Lowering Stack} \\[0.6em]
  {\normalsize \textit{PhaseMirror-HQ / Multiplicity Foundation}} \\[0.3em]
  {\normalsize Repository: \texttt{github.com/MultiplicityFoundation/PhaseMirror-HQ}} \\[0.3em]
  {\normalsize Publication Date: April 25, 2026 \quad Version: Gate U\ensuremath{^{\star}}}
}
\author{}
\date{}

%% ─── Document ─────────────────────────────────────────────────────────────────
\begin{document}

\maketitle
\thispagestyle{fancy}

\begin{defensebox}
This document is published as a \textbf{prior art defensive publication} under the
Multiplicity Foundation's open-IP policy. It discloses, with full technical
specificity, the design of Gate U: a hardware acceleration roadmap for the
Prime-Indexed Recursive Tensor Machine (PIRTM) dialect. Publication of this
document establishes prior art as of April 25, 2026, preventing downstream
patent capture of the disclosed techniques without attribution to this work.
All specifications herein are normative for the PhaseMirror-HQ codebase.
\end{defensebox}

\vspace{0.5em}

%% ─── Abstract ────────────────────────────────────────────────────────────────
\begin{abstract}
\noindent
Gate U is the point in the Multiplicity Theory roadmap where the PIRTM MLIR
dialect transitions from a software abstraction to a hardware specification.
This document presents a comprehensive technical overview of Gate U and its
successor specification Gate U\ensuremath{^{\star}}, covering: (1) the mathematical basis
for prime-indexed tensor contraction as a hardware datapath; (2) the three-tier
hardware feasibility framework (Tier 0 simulation, Tier 1 FPGA emulation,
Tier 2 ASIC prototype); (3) the MLIR TableGen and lowering pass for the
\texttt{pirtm.step.hw} operation; (4) an HLS kernel for FPGA prototype
implementation; (5) the hardware constants charter (ADR-037) declaring all
normative constants including the hardware convergence threshold
\(\varepsilon_{\mathrm{hw}} = 2^{-20}\), contraction ceiling
\(r_{\mathrm{hw}} = 0.9\), and canonical spectral pipeline depth
\(d_{\mathrm{hw}} = 132\); (6) the irreducibility certificate contract as a
first-class artifact; (7) governance causality locking via the digital twin
commitment hash; and (8) the three-verdict spectral pipeline
(\texttt{CONTRACTIVE}, \texttt{VIOLATION}, \texttt{SPECTRAL\_UNCERTAIN}).
This publication establishes prior art for hardware-aware governance enforcement
at the silicon layer in distributed prime-indexed computation systems.
\end{abstract}

\tableofcontents
\newpage

%% ═══════════════════════════════════════════════════════════════════════════
\section{Executive Summary}
%% ═══════════════════════════════════════════════════════════════════════════

Gate U marks the first point in the PhaseMirror roadmap where Multiplicity
Theory's algebraic identity structure---prime-indexed session identity---becomes
the organizing principle of a \emph{physical hardware datapath}, not merely a
software label. The core claim: a prime index \(p_i\) is not metadata attached
to a computation; it \emph{is} the routing key that shapes silicon interconnect,
fixes spectral pipeline depth, and constrains the weight matrix structure on-die.

\subsection{Key Contributions}

\begin{enumerate}[label=\textbf{C\arabic*.}]
  \item \textbf{Prime-routed hardware datapath.} The \texttt{pirtm.step}
        operator---a contractive map on a prime-indexed multiplicity
        space---maps directly to a prime-labeled routing fabric. This is the
        first disclosed design where primality is the hardware routing predicate
        at the ISA level.

  \item \textbf{Three-tier feasibility framework.} ADR-SHF-006 defines Tier 0
        (simulation), Tier 1 (FPGA emulation with hardware-aware constraints),
        and Tier 2 (physical ASIC prototype). Gate U traverses all three tiers
        with explicit evidence requirements and governance sign-off.

  \item \textbf{Hardware constants charter (ADR-037).} Five hardware constants
        are declared normatively, eliminating all magic numbers from the
        implementation stack. The canonical spectral pipeline depth
        \(d_{\mathrm{hw}} = 132\) is derived from
        \(\varepsilon_{\mathrm{hw}}\) and \(r_{\mathrm{hw}}\), not hardcoded.

  \item \textbf{Three-verdict spectral pipeline.} The spectral pipeline emits
        one of three verdicts---\texttt{CONTRACTIVE}, \texttt{VIOLATION},
        \texttt{SPECTRAL\_UNCERTAIN}---where the third verdict detects
        degenerate eigenvalue gaps and routes to software fallback, closing
        the false-\texttt{CONTRACTIVE} safety hole.

  \item \textbf{Irreducibility certificate as first-class artifact.} Primality
        alone does not guarantee irreducibility of the weight matrix
        representation. An irreducibility certificate (rank check on
        \(W^{(p_i)} - \hat{\rho} I\)) is required before ASIC-specific
        routing optimizations are applied.

  \item \textbf{Governance causality lock.} Tier promotion requires a
        cryptographic commitment hash \(C = H(\text{tier} \| t \| p_i \|
        \varepsilon_{\mathrm{hw}} \| r_{\mathrm{hw}} \| d_{\mathrm{hw}})\)
        embedded in the digital twin snapshot and verified by the hardware
        driver before dispatch. No hardware activation precedes governance sign-off.

  \item \textbf{Silicon-layer Byzantine fault identification.} A false
        \texttt{CONTRACTIVE} verdict propagating in a governance mesh constitutes
        a Byzantine fault injected at the silicon layer, bypassing all
        software-layer governance checks. This failure mode is named and
        bounded in ADR-037's ADR-038 stub.
\end{enumerate}

\subsection{Performance Targets}

On a Xilinx UltraScale+ \texttt{xcvu9p}-class FPGA at Tier 1:
\begin{itemize}
  \item Throughput: $\approx$333,000 \texttt{pirtm.step} operations/second at 300~MHz
  \item Pipeline latency: $\approx$840 cycles (corrected from 900 with \(d_{\mathrm{hw}}=132\))
  \item Spectral check: included inline, contractivity gate hardware-enforced
  \item Gain over software: $\approx$6$\times$ additional speedup over the Tier 0
        software target of 50K steps/sec on modern CPU
\end{itemize}

At Tier 2 (TSMC 5nm representative node), total chip power is estimated at
$\approx$257~mW with the governance register always-on and
non-power-gateable.

%% ═══════════════════════════════════════════════════════════════════════════
\section{Mathematical Foundations}
%% ═══════════════════════════════════════════════════════════════════════════

\subsection{Multiplicity Spaces and Prime-Indexed Dynamics}

\begin{definition}[Prime Index Domain]
Let \(\mathcal{P}_{\max}\) denote the declared prime index domain:
\[
  \mathcal{P}_{\max} = \bigl\{ p \in \mathbb{Z}^+ : p \text{ prime},\; p \leq P_{\max} \bigr\},
  \quad P_{\max} = 2^{31} - 1 = 2{,}147{,}483{,}647.
\]
\(P_{\max}\) is the largest 32-bit prime. All session identifiers \(p_i\) must
satisfy \(p_i \in \mathcal{P}_{\max}\).
\end{definition}

\begin{definition}[Multiplicity Space]
For \(p \in \mathcal{P}_{\max}\), the prime-indexed multiplicity space is
\(\mathcal{V}_p = \mathbb{R}^n\) equipped with the Euclidean norm
\(\|\cdot\|_2\). For the Gate U reference implementation, \(n = 512\).
\end{definition}

\begin{definition}[Prime-Indexed Contractive Map]
The \texttt{pirtm.step} operation is the linear map:
\[
  T_p : \mathcal{V}_p \to \mathcal{V}_p, \qquad T_p(x) = W^{(p)} x,
\]
where \(W^{(p)} \in \mathbb{R}^{n \times n}\) is the session weight matrix and
\(\rho(W^{(p)}) < r_{\mathrm{hw}}\) is the hardware contractivity invariant.
\end{definition}

\subsubsection{Structural Sparsity from Prime Factorization}

The nonzero pattern of \(W^{(p)}\) is constrained by the divisibility lattice
of \(p\). Because \(p\) is prime, it has no proper divisors; its divisibility
lattice is trivial \(\{1, p\}\). Therefore \(W^{(p)}\) operates in a single
irreducible subspace---the weight matrix is \emph{structurally} block-diagonal
with one block. This sparsity is not accidental: it is the algebraic consequence
of prime identity and can be hardwired in silicon.

\begin{remark}
Structural sparsity from prime factorization is the key distinction between a
PIRTM step engine and a general matrix accelerator (e.g., a TPU systolic array).
The silicon datapath shape encodes the mathematical identity of the session.
\end{remark}

\subsection{Spectral Radius and Contractivity}

\begin{definition}[Spectral Radius]
\[
  \rho(W) = \lim_{n \to \infty} \|W^n\|^{1/n} = \max_i |\lambda_i(W)|.
\]
\end{definition}

\begin{theorem}[Geometric Convergence Bound]
If \(\rho(W) \leq r < 1\), then there exists \(C > 0\) such that
\[
  \|W^n\| \leq C \cdot r^n \quad \forall\, n \geq 0.
\]
The convergence rate \(r\) determines the required spectral pipeline depth.
\end{theorem}

\subsubsection{Hardware Convergence Threshold and Canonical Depth}

\begin{invariant}[Hardware Constants Charter, ADR-037]\label{inv:constants}
The following five constants are declared normatively. All code-level uses
\textbf{must} derive from these constants by import or code generation.
Magic number overrides are forbidden in hardware-targeting code paths.

\begin{center}
\begin{tabular}{lllp{5cm}}
\toprule
\textbf{Identifier} & \textbf{Symbol} & \textbf{Value} & \textbf{Derivation} \\
\midrule
\texttt{HW\_EPSILON}             & \(\varepsilon_{\mathrm{hw}}\) & \(2^{-20}\)        & Conservative float32 convergence threshold \\
\texttt{HW\_CONTRACTION\_CEILING}& \(r_{\mathrm{hw}}\)           & \(0.9\)            & Max bound for Tier 1 timing closure \\
\texttt{HW\_SPECTRAL\_DEPTH}     & \(d_{\mathrm{hw}}\)           & \(132\)            & \(\lceil \log(\varepsilon_{\mathrm{hw}}) / \log(r_{\mathrm{hw}}) \rceil\) \\
\texttt{HW\_PRIME\_MAX}          & \(P_{\max}\)                  & \(2{,}147{,}483{,}647\) & Largest 32-bit prime \\
\texttt{HW\_EIGENVALUE\_GAP}     & \(\delta_{\mathrm{gap}}\)     & \(0.01\)           & Gap below which \texttt{SPECTRAL\_UNCERTAIN} fires \\
\bottomrule
\end{tabular}
\end{center}
\end{invariant}

\begin{proposition}[Canonical Spectral Depth Derivation]
The canonical pipeline depth is:
\begin{equation}
  d_{\mathrm{hw}} = \left\lceil \frac{\log \varepsilon_{\mathrm{hw}}}{\log r_{\mathrm{hw}}} \right\rceil
  = \left\lceil \frac{-20 \ln 2}{\ln 0.9} \right\rceil
  = \left\lceil \frac{-13.8629}{-0.10536} \right\rceil
  = \lceil 131.56 \rceil = 132. \label{eq:depth}
\end{equation}
\end{proposition}

\begin{remark}
The value 150 previously appearing in \texttt{hw\_dispatch.py} and the HLS
kernel is incorrect. It corresponds neither to \(\varepsilon_{\mathrm{hw}} = 10^{-6}\)
at \(r = 0.9\) (which gives 132) nor to float32 machine epsilon at \(r = 0.9\)
(which gives 152). It is a magic number with no derivation. ADR-037 replaces
it with \(d_{\mathrm{hw}} = 132\) from \cref{eq:depth}.
\end{remark}

\subsection{Three-Verdict Spectral Pipeline}

Let \(\hat{\rho}_d\) be the power iteration estimate after \(d_{\mathrm{hw}}\)
steps, and let
\[
  \Delta\lambda = \frac{|\lambda_1|}{|\lambda_2|} - 1
\]
be the relative eigenvalue gap (estimated via deflation).

\begin{definition}[Spectral Verdict Function]
\begin{equation}
  \Sigma(W^{(p)}) = \begin{cases}
    \texttt{CONTRACTIVE}        & \hat{\rho}_d < r_{\mathrm{hw}} \text{ and } \Delta\lambda \geq \delta_{\mathrm{gap}} \\
    \texttt{VIOLATION}          & \hat{\rho}_d \geq r_{\mathrm{hw}} \\
    \texttt{SPECTRAL\_UNCERTAIN} & \hat{\rho}_d < r_{\mathrm{hw}} \text{ and } \Delta\lambda < \delta_{\mathrm{gap}}
  \end{cases} \label{eq:verdicts}
\end{equation}
\texttt{SPECTRAL\_UNCERTAIN} routes to Tier 0 software fallback. It is a
\emph{precision signal}, not a governance violation.
\end{definition}

\begin{theorem}[Correctness Under Gap Condition]
If \(\Delta\lambda \geq \delta_{\mathrm{gap}}\), then after \(d_{\mathrm{hw}}\)
power iterations:
\[
  \bigl|\hat{\rho}_{d_{\mathrm{hw}}} - \rho(W^{(p)})\bigr|
  \leq \varepsilon_{\mathrm{hw}} \cdot \rho(W^{(p)}).
\]
Therefore the \texttt{CONTRACTIVE} verdict is correct whenever
\(\rho(W^{(p)}) < r_{\mathrm{hw}} - \varepsilon_{\mathrm{hw}}\).
\end{theorem}

\subsection{Irreducibility Certificate}

\begin{definition}[Irreducibility Certificate]
A weight matrix \(W^{(p)}\) holds an \textbf{irreducibility certificate} if and
only if:
\begin{equation}
  \mathrm{rank}\bigl(W^{(p)} - \hat{\rho}\, I_n\bigr) = n - 1, \label{eq:irred}
\end{equation}
where \(\hat{\rho} = \hat{\rho}_{d_{\mathrm{hw}}}\) is the power iteration
estimate and \(n\) is the matrix dimension.
\end{definition}

\begin{remark}
This rank check is \(\mathcal{O}(n^2)\)---cheaper than a full
eigendecomposition---and is implementable as the Tier 0 reference without any
hardware dependency. If \cref{eq:irred} fails (rank \(> n-1\)), the eigenspace
is degenerate and the session routes to the general crossbar path; ASIC-specific
block-diagonal routing optimizations must not be applied.
\end{remark}

\begin{contract}[ASIC Routing Precondition]
ASIC-specific routing optimizations (single-block-diagonal weight layout,
no inter-session crossbar) \textbf{must not} be applied unless
\(\mathrm{Irred}(W^{(p)}) = \mathrm{True}\). The certificate artifact is the
tuple \((p_i, \hat{\rho}, \mathrm{rank\_result}, t)\) stored in the digital
twin snapshot alongside the primality certificate.
\end{contract}

\subsection{Tier Promotion Predicate}

Let \(\mathcal{E}_{\min}(k)\) be the minimum evidence closure set required for
hardware tier \(k\), and let \(\sigma\) be the safety checklist signature. The
promotion is valid if and only if:

\begin{equation}
  \Pi(\Gamma) = 1 \iff
  \mathcal{E} \supseteq \mathcal{E}_{\min}(\mathrm{tier}_{\mathrm{target}})
  \;\wedge\;
  \bigl(\mathrm{tier}_{\mathrm{target}} < 2 \;\vee\; \sigma \neq \emptyset\bigr). \label{eq:promotion}
\end{equation}

\subsection{Governance Commitment Hash}

\begin{definition}[Tier Commitment Hash]
Let \(\mathcal{K} = (\varepsilon_{\mathrm{hw}}, r_{\mathrm{hw}}, d_{\mathrm{hw}})\)
be the ADR-037 hardware constant triple. The commitment hash is:
\begin{equation}
  C = H\bigl(\mathrm{tier} \;\|\; t \;\|\; p_i \;\|\; \varepsilon_{\mathrm{hw}} \;\|\; r_{\mathrm{hw}} \;\|\; d_{\mathrm{hw}}\bigr), \label{eq:hash}
\end{equation}
where \(H\) is a collision-resistant hash function (e.g., SHA-256).
\end{definition}

\begin{remark}
Including \(\mathcal{K}\) in the commitment ensures that any bitstream compiled
against stale hardware constants (e.g., old \(d_{\mathrm{hw}} = 150\)) produces
a different \(C\), which the hardware driver rejects before dispatch. This makes
ADR-037 constant drift a commitment verification failure detectable at driver
handshake time.
\end{remark}

%% ═══════════════════════════════════════════════════════════════════════════
\section{Hardware Feasibility Tiers (ADR-SHF-006)}
%% ═══════════════════════════════════════════════════════════════════════════

ADR-SHF-006 defines three feasibility tiers that Gate U must traverse in
sequence. Every hardware claim must carry a tier label, assumptions, and
limitations.

\begin{table}[H]
\centering
\caption{Hardware Feasibility Tier Framework (ADR-SHF-006)}
\label{tab:tiers}
\begin{tabular}{cllp{4.5cm}}
\toprule
\textbf{Tier} & \textbf{Label} & \textbf{Gate U Milestone} & \textbf{Evidence Required} \\
\midrule
0 & Simulated only      & PIRTM Python/NumPy simulation        & Deterministic reproducer, fixed seed (ADR-SHF-002) \\
1 & Emulated, HW-aware  & FPGA prototype via MLIR$\to$HLS      & Timing report, resource utilization, spectral pass rate \\
2 & Physical prototype  & ASIC tape-out or hard IP block       & Silicon measurement, power budget, error rate \\
\bottomrule
\end{tabular}
\end{table}

Claim \textbf{SHF-C012} in the traceability matrix---``hardware pathways are
feasible across quantum, neuromorphic, and hybrid stacks''---is currently open
with required evidence tiers S and E. Gate U closes it. The tier promotion
predicate from \cref{eq:promotion} governs all tier transitions.

%% ═══════════════════════════════════════════════════════════════════════════
\section{MLIR Dialect Lowering Architecture}
%% ═══════════════════════════════════════════════════════════════════════════

\subsection{Lowering Chain}

The PIRTM MLIR dialect defines \texttt{pirtm.step} as a high-level tensor
contraction operation with spectral radius verification. The Gate U lowering
chain is:

\begin{center}
\begin{tabular}{c}
\texttt{Python (pirtm\_core/)} \\
$\downarrow$ \\
\texttt{MLIR PIRTM dialect (pirtm.step, pirtm.reason)} \\
$\downarrow$ \\
\texttt{MLIR Linalg dialect (generic tensor ops)} \\
$\downarrow$ \\
\texttt{MLIR affine + memref dialects} \\
$\swarrow$ \hspace{3em} $\searrow$ \\
\texttt{CPU: LLVM IR $\to$ binary} \hspace{1em} \texttt{HW: HLS C++ $\to$ RTL $\to$ FPGA/ASIC}
\end{tabular}
\end{center}

The Gate U contribution is the \emph{hardware lowering pass}:
a new MLIR pass converting \texttt{pirtm.step} to HLS-compatible IR with
prime-indexed routing constraints explicit in the representation.

\subsection{TableGen: \texttt{pirtm.step.hw} Operation}

\begin{lstlisting}[style=tablegen, caption={\texttt{PIRTMOps.td} --- Gate U hardware op definition}, label=lst:tablegen]
// pirtm/Dialect/PIRTM/IR/PIRTMOps.td — Gate U additions

def PIRTM_StepHWOp : PIRTM_Op<"step.hw", [
    Pure,
    DeclareOpInterfaceMethods<ContractionVerifiable>,
    DeclareOpInterfaceMethods<HardwareMappable>
]> {
    let summary = "Hardware-lowered pirtm.step for FPGA/ASIC targets";
    let description = [{
        Lowers a pirtm.step to a hardware-aware representation.
        prime_id drives index routing in the generated RTL.
        spectral_pipeline_depth is derived from HW_SPECTRAL_DEPTH (ADR-037):
          depth = ceil(log(eps_hw) / log(r_hw)) = 132
        contraction_bound must satisfy: bound <= HW_CONTRACTION_CEILING = 0.9
    }];
    let arguments = (ins
        MemRefType:$input_buf,           // DMA-mapped input tensor
        MemRefType:$weight_buf,          // weight matrix W^(p_i)
        MemRefType:$output_buf,          // DMA-mapped output tensor
        I64Attr:$prime_id,               // routing key; must pass Miller-Rabin
        F32Attr:$contraction_bound,      // hardware enforces rho < bound
        I32Attr:$spectral_pipeline_depth,// fixed pipeline depth (= 132)
        PIRTM_HWTierAttr:$hw_tier        // TIER0, TIER1, TIER2
    );
    let results = (outs
        I2:$spectral_verdict,  // 0=CONTRACTIVE, 1=VIOLATION, 2=UNCERTAIN
        I64:$cycle_count       // profiling: cycles consumed
    );
    let verifier = [{ return verifyHardwareTierConstraints(*this); }];
    let hasCustomAssemblyFormat = 1;
}

// Hardware tier enum
def PIRTM_HWTierAttr : I32EnumAttr<"HWTier", "Hardware feasibility tier", [
    I32EnumAttrCase<"TIER0", 0>,   // simulated
    I32EnumAttrCase<"TIER1", 1>,   // FPGA emulated
    I32EnumAttrCase<"TIER2", 2>    // ASIC prototype
]>;

// Spectral verdict enum (three-verdict pipeline, ADR-037)
def PIRTM_SpectralVerdictAttr : I32EnumAttr<"SpectralVerdict",
    "Three-verdict spectral pipeline output", [
    I32EnumAttrCase<"CONTRACTIVE",        0>,
    I32EnumAttrCase<"VIOLATION",          1>,
    I32EnumAttrCase<"SPECTRAL_UNCERTAIN", 2>
]>;
\end{lstlisting}

\subsection{MLIR Lowering Pass (C++)}

\begin{lstlisting}[style=cppstyle, caption={\texttt{PIRTMToHWPass.cpp} --- hardware lowering pass}, label=lst:pass]
// lib/Conversion/PIRTMToHW/PIRTMToHWPass.cpp

#include "pirtm/HW/PIRTMHWConstants.h"  // single source: ADR-037 constants

struct PIRTMStepToHWPattern : public OpRewritePattern<pirtm::StepOp> {
    using OpRewritePattern::OpRewritePattern;

    LogicalResult matchAndRewrite(pirtm::StepOp op,
                                  PatternRewriter &rewriter) const override {
        // 1. Validate prime_id via Miller-Rabin (20 rounds)
        int64_t primeId = op.getPrimeId();
        if (primeId > PIRTM_HW_PRIME_MAX) {
            return op.emitError("prime_id exceeds P_max=2^31-1 (ADR-037 L0.5)");
        }
        if (!millerRabinCheck(primeId, /*rounds=*/20)) {
            return op.emitError("prime_id fails primality; L0 invariant violated");
        }

        // 2. Enforce hardware contraction bound <= r_hw (ADR-037)
        float bound = op.getContractionBound();
        if (bound > PIRTM_HW_CONTRACTION_CEILING) {
            return op.emitError("contraction_bound ")
                << bound << " > r_hw=" << PIRTM_HW_CONTRACTION_CEILING
                << "; use HardwareContractionBound type (ADR-037)";
        }

        // 3. Use canonical spectral depth from ADR-037, not op attribute
        int depth = PIRTM_HW_SPECTRAL_DEPTH;  // = 132, never override

        // 4. Emit step.hw op with routing metadata
        auto hwOp = rewriter.create<pirtm::StepHWOp>(
            op.getLoc(),
            op.getInputBuf(), op.getWeightBuf(), op.getOutputBuf(),
            rewriter.getI64IntegerAttr(primeId),
            rewriter.getF32FloatAttr(bound),
            rewriter.getI32IntegerAttr(depth),
            pirtm::HWTierAttr::get(op.getContext(), pirtm::HWTier::TIER1)
        );
        rewriter.replaceOp(op, hwOp.getResults());
        return success();
    }
};
\end{lstlisting}

%% ═══════════════════════════════════════════════════════════════════════════
\section{FPGA Prototype --- Tier 1}
%% ═══════════════════════════════════════════════════════════════════════════

\subsection{HLS Kernel: \texttt{pirtm\_step}}

The hardware lowering pass emits HLS C++ synthesized by Vitis HLS or Intel
oneAPI HLS to RTL. The prime index drives the routing selector; the spectral
pipeline runs as a parallel accumulator pipeline.

\begin{lstlisting}[style=cppstyle, caption={HLS kernel \texttt{pirtm\_step\_kernel.cpp}}, label=lst:hls]
// hls/pirtm_step_kernel.cpp — Tier 1 FPGA prototype

#include <ap_fixed.h>
#include <hls_stream.h>
#include "pirtm_hw_constants.h"  // ADR-037: PIRTM_HW_SPECTRAL_DEPTH=132, etc.

typedef ap_fixed<32, 8, AP_RND, AP_SAT> pirtm_fp;  // Q8.24 fixed-point
typedef ap_uint<64> prime_id_t;

// Spectral verdict enum matching MLIR SpectralVerdictAttr
enum SpectralVerdict { CONTRACTIVE = 0, VIOLATION = 1, SPECTRAL_UNCERTAIN = 2 };

// Deflation gap estimation: compare dominant vs. subdominant eigenvalue
pirtm_fp estimate_gap(pirtm_fp W[512][512], pirtm_fp v[512], int dim) {
#pragma HLS INLINE off
    // Deflate: compute Av - rho*v, measure residual as gap proxy
    pirtm_fp rho_v[512], residual = 0;
    for (int j = 0; j < dim; j++) {
        pirtm_fp wv = 0;
        for (int k = 0; k < dim; k++) wv += W[j][k] * v[k];
        rho_v[j] = wv;
    }
    // Gap proxy: ||Av||/||v|| ratio for second vector (simplified)
    pirtm_fp norm1 = 0, norm2 = 0;
    for (int j = 0; j < dim; j++) { norm1 += v[j]*v[j]; norm2 += rho_v[j]*rho_v[j]; }
    return hls::sqrt(norm2) / (hls::sqrt(norm1) + 1e-9f);
}

SpectralVerdict spectral_pipeline(
    pirtm_fp W[512][512], int dim,
    pirtm_fp contraction_bound, pirtm_fp* rho_out
) {
#pragma HLS INLINE off
    pirtm_fp v[512], v_new[512], norm = 0;
    for (int i = 0; i < dim; i++) v[i] = (i == 0) ? 1.0 : 0.0;

    for (int iter = 0; iter < PIRTM_HW_SPECTRAL_DEPTH; iter++) {
#pragma HLS PIPELINE II=1
        norm = 0;
        for (int j = 0; j < dim; j++) {
#pragma HLS UNROLL factor=8
            v_new[j] = 0;
            for (int k = 0; k < dim; k++) v_new[j] += W[j][k] * v[k];
            norm += v_new[j] * v_new[j];
        }
        norm = hls::sqrt(norm);
        for (int j = 0; j < dim; j++) v[j] = v_new[j] / norm;
    }
    *rho_out = norm;

    if (norm >= contraction_bound) return VIOLATION;

    // Eigenvalue gap check (ADR-037: delta_gap = PIRTM_HW_EIGENVALUE_GAP)
    pirtm_fp gap = estimate_gap(W, v, dim) / norm - 1.0f;
    if (gap < PIRTM_HW_EIGENVALUE_GAP) return SPECTRAL_UNCERTAIN;

    return CONTRACTIVE;
}

void pirtm_step(
    pirtm_fp input[512], pirtm_fp W[512][512], pirtm_fp output[512],
    prime_id_t prime_id, pirtm_fp contraction_bound,
    ap_uint<2>& spectral_verdict,   // 0=CONTRACTIVE,1=VIOLATION,2=UNCERTAIN
    ap_uint<32>& cycle_count
) {
#pragma HLS INTERFACE m_axi port=input  bundle=gmem0
#pragma HLS INTERFACE m_axi port=W      bundle=gmem1
#pragma HLS INTERFACE m_axi port=output bundle=gmem2
#pragma HLS INTERFACE s_axilite port=prime_id
#pragma HLS INTERFACE s_axilite port=spectral_verdict
#pragma HLS INTERFACE s_axilite port=return

    int dim = 512;
    pirtm_fp rho;
    SpectralVerdict sv = spectral_pipeline(W, dim, contraction_bound, &rho);
    spectral_verdict = (ap_uint<2>)sv;

    if (sv == CONTRACTIVE) {
        for (int j = 0; j < dim; j++) {
#pragma HLS PIPELINE II=1
            pirtm_fp acc = 0;
            for (int k = 0; k < dim; k++) {
#pragma HLS UNROLL factor=16
                acc += W[j][k] * input[k];
            }
            output[j] = acc;
        }
    }

    // Governance: VIOLATION writes non-maskable interrupt register
    // AXI4-Lite PIRTM_GOVERNANCE_REG bit 0 = VIOLATION flag
    // ADR-031 daemon watchdog reads on heartbeat; ADR-028 kill-switch escalates
}
\end{lstlisting}

\subsection{FPGA Resource Utilization Estimate}

\begin{table}[H]
\centering
\caption{Tier 1 resource utilization on Xilinx UltraScale+ \texttt{xcvu9p}}
\label{tab:resources}
\begin{tabular}{lrrl}
\toprule
\textbf{Resource} & \textbf{Estimated} & \textbf{Budget} & \textbf{Notes} \\
\midrule
DSP slices (tensor MAC)     & $\sim$3,200 & 6,840    & 512 $\times$ 8-wide SIMD \\
BRAM 18K (weight matrix)    & 128         & 4,320    & 512$\times$512 $\times$ 4B \\
LUT (routing + control)     & $\sim$45K   & 1.18M    & Prime router + state \\
Pipeline latency            & $\approx$840 cycles & --- & At 300 MHz \\
Throughput                  & $\approx$333K steps/sec & --- & Spectral check included \\
\bottomrule
\end{tabular}
\end{table}

\begin{remark}
The latency of 840 cycles represents an 8\% improvement over the initial 900-cycle
estimate, achieved by correcting \(d_{\mathrm{hw}} = 132\) from the erroneous
value of 150.
\end{remark}

%% ═══════════════════════════════════════════════════════════════════════════
\section{ASIC Design Exploration --- Tier 2}
%% ═══════════════════════════════════════════════════════════════════════════

\subsection{Prime-Routing Silicon Architecture}

A PIRTM ASIC is not a general matrix accelerator. It is a prime-indexed tensor
engine whose datapath structure reflects the mathematical identity of each session.
Key design insight from ADR-SHF-002: prime-indexed dynamics operate in irreducible
subspaces, so \(W^{(p_i)}\) for a prime session is block-diagonal with one block.
The ASIC hardwires this---no general crossbar is needed for intra-session compute.
The crossbar appears only at the inter-session composition boundary
(\(p_i \cdot p_j\)).

\begin{figure}[H]
\centering
\begin{minipage}{0.85\textwidth}
\begin{verbatim}
┌─────────────────────────────────────────────────────────┐
│              PIRTM Step Engine (ASIC die)                │
│                                                          │
│  ┌─────────────┐  prime_id  ┌──────────────────────┐   │
│  │ Prime Router│───────────►│  Session Dispatch Unit │   │
│  │ (factorizer)│            │  maps p_i → lane[i]   │   │
│  └─────────────┘            └──────────┬───────────┘   │
│                                         │                │
│  ┌──────────────┐  ┌──────────────┐  ┌─────────────┐  │
│  │  MAC Array   │  │  Spectral    │  │  Governance │  │
│  │ (tensor      │  │  Pipeline    │  │  Register   │  │
│  │  contraction)│  │  power iter  │  │ (AXI4-Lite) │  │
│  └──────┬───────┘  └──────┬───────┘  └──────┬──────┘  │
│         ▼                 ▼                  ▼           │
│    output_buf      spectral_verdict    VIOLATION_IRQ    │
└─────────────────────────────────────────────────────────┘
\end{verbatim}
\end{minipage}
\caption{PIRTM Step Engine ASIC die block diagram}
\label{fig:asic}
\end{figure}

The \texttt{VIOLATION\_IRQ} line is hardwired to a non-maskable interrupt (NMI).
No software path can suppress it. This is the silicon expression of ADR-028's
kill-switch non-negotiable.

\subsection{Governance Register SystemVerilog Stub}

\begin{lstlisting}[style=verilogstyle, caption={Governance register NMI stub (ADR-028/ADR-032)}, label=lst:sv]
// rtl/governance_reg.sv  — Tier 2 ASIC governance stub

module governance_reg #(
    parameter BASE_ADDR = 32'hFF00_0000
) (
    input  logic clk,
    input  logic rst_n,
    input  logic violation_in,    // from spectral pipeline
    input  logic uncertain_in,    // SPECTRAL_UNCERTAIN: log but no NMI
    output logic nmi_out,         // non-maskable interrupt to CPU
    output logic uncertain_flag,  // advisory flag to daemon watchdog
    // AXI4-Lite slave interface (read-only; no software clear path)
    input  logic [31:0] s_axilite_araddr,
    output logic [31:0] s_axilite_rdata,
    output logic        s_axilite_rvalid
);
    // NMI: non-maskable, no software clear, always-on power domain
    always_ff @(posedge clk or negedge rst_n)
        nmi_out <= rst_n ? violation_in : 1'b0;

    // SPECTRAL_UNCERTAIN: advisory, does not trigger NMI
    always_ff @(posedge clk or negedge rst_n)
        uncertain_flag <= rst_n ? uncertain_in : 1'b0;

    // AXI4-Lite read: expose both flags to software for logging
    always_ff @(posedge clk) begin
        if (s_axilite_araddr == BASE_ADDR) begin
            s_axilite_rdata  <= {30'b0, uncertain_flag, nmi_out};
            s_axilite_rvalid <= 1'b1;
        end else
            s_axilite_rvalid <= 1'b0;
    end
endmodule
\end{lstlisting}

\subsection{ASIC Power Budget}

\begin{table}[H]
\centering
\caption{Tier 2 power budget (TSMC 5nm, representative 2032 node, 1~GHz)}
\label{tab:power}
\begin{tabular}{lrl}
\toprule
\textbf{Block} & \textbf{Power (mW)} & \textbf{Notes} \\
\midrule
MAC Array (512 FP32 MACs)    & $\sim$180 & Pipelined, 1 GHz target \\
Spectral Pipeline             & $\sim$40  & Fixed depth 132 cycles \\
Prime Router                  & $\sim$5   & Factorization LUT, fixed size \\
Governance Register + NMI IRQ & $\sim$2   & Always-on, non-power-gateable \\
DMA / AXI Interconnect        & $\sim$30  & Bandwidth-limited \\
\midrule
\textbf{Total}               & $\sim$\textbf{257} & Sub-300~mW budget met \\
\bottomrule
\end{tabular}
\end{table}

The governance register is always-on---it cannot be power-gated. This is an
explicit safety constraint from ADR-SHF-006.

%% ═══════════════════════════════════════════════════════════════════════════
\section{Python Runtime: Hardware Dispatch Layer}
%% ═══════════════════════════════════════════════════════════════════════════

\subsection{Hardware Constants Module (Single Source of Truth)}

\begin{lstlisting}[style=pythonstyle, caption={\texttt{pirtm\_core/hw\_constants.py} --- ADR-037 normative constants}, label=lst:constants]
# pirtm_core/hw_constants.py  — ADR-037: Hardware Constants Charter
# Single source of truth. All hardware-targeting code imports from here.
# NEVER hardcode these values anywhere else in the codebase.

import math
import numpy as np

# --- ADR-037 Normative Constants -------------------------------------------
HW_EPSILON: float           = 2.0 ** -20         # 9.537e-7
HW_CONTRACTION_CEILING: float = 0.9
HW_SPECTRAL_DEPTH: int      = math.ceil(
    math.log(HW_EPSILON) / math.log(HW_CONTRACTION_CEILING)
)  # = 132; computed, never hardcoded
HW_PRIME_MAX: int           = 2_147_483_647       # 2^31 - 1, largest 32-bit prime
HW_EIGENVALUE_GAP: float    = 0.01

assert HW_SPECTRAL_DEPTH == 132, \
    f"Depth invariant broken: expected 132, got {HW_SPECTRAL_DEPTH}"

# --- Irreducibility Certificate ---------------------------------------------
def irreducibility_check(W: np.ndarray, rho_hat: float) -> bool:
    """
    Returns True iff W has a single dominant eigenspace (irreducible).
    Contract (ADR-037): must return True before ASIC block-diagonal
    routing optimizations are applied.
    Complexity: O(n^2) via matrix_rank on (W - rho_hat * I).
    """
    n = W.shape[0]
    rank = np.linalg.matrix_rank(
        W - rho_hat * np.eye(n), tol=HW_EPSILON
    )
    return rank == n - 1

# --- HardwareContractionBound type -----------------------------------------
class HardwareContractionBound:
    """
    Type-safe hardware contraction bound.
    Enforces bound <= HW_CONTRACTION_CEILING at construction time.
    Use this type for all Tier 1 and Tier 2 dispatch configs.
    """
    def __init__(self, value: float):
        if value > HW_CONTRACTION_CEILING:
            raise ValueError(
                f"Hardware bound {value} > r_hw={HW_CONTRACTION_CEILING}. "
                f"Use SoftwareContractionBound for Tier 0."
            )
        if value <= 0.0 or value >= 1.0:
            raise ValueError(f"Bound must be in (0, 1), got {value}")
        self.value = value

class SoftwareContractionBound:
    """Tier 0 only. No hardware constraint on value."""
    def __init__(self, value: float):
        assert 0.0 < value < 1.0
        self.value = value
\end{lstlisting}

\subsection{Hardware Dispatch Orchestrator}

\begin{lstlisting}[style=pythonstyle, caption={\texttt{pirtm\_core/hw\_dispatch.py} --- Gate U hardware dispatch layer}, label=lst:dispatch]
# pirtm_core/hw_dispatch.py  — Gate U hardware dispatch layer

import hashlib, struct, ctypes, os
import numpy as np
from enum import Enum
from dataclasses import dataclass
from pirtm_core.hw_constants import (
    HW_EPSILON, HW_CONTRACTION_CEILING, HW_SPECTRAL_DEPTH,
    HW_PRIME_MAX, HW_EIGENVALUE_GAP,
    HardwareContractionBound, SoftwareContractionBound,
    irreducibility_check
)

class SpectralVerdict(Enum):
    CONTRACTIVE        = 0
    VIOLATION          = 1
    SPECTRAL_UNCERTAIN = 2

class HWTier(Enum):
    TIER0 = 0  # Python/NumPy simulation
    TIER1 = 1  # FPGA via XRT (Xilinx Runtime)
    TIER2 = 2  # ASIC via PCIe driver

@dataclass
class HWDispatchConfig:
    tier: HWTier
    device_path: str
    prime_id: int
    contraction_bound: HardwareContractionBound  # type-safe; rejects > 0.9
    tier_commitment_hash: str = ""  # must be set before Tier 1/2 dispatch

    def __post_init__(self):
        if self.prime_id > HW_PRIME_MAX:
            raise ValueError(f"prime_id {self.prime_id} > P_max (ADR-037)")
        if self.tier != HWTier.TIER0 and not self.tier_commitment_hash:
            raise RuntimeError(
                "Tier commitment hash required for Tier 1/2 dispatch (ADR-032)"
            )

def compute_tier_commitment(tier: int, prime_id: int, timestamp: float) -> str:
    """
    Commitment hash per ADR-032 amendment.
    Includes all ADR-037 constants to detect stale-constant bitstreams.
    C = H(tier || t || p_i || eps_hw || r_hw || d_hw)
    """
    payload = struct.pack(
        ">iqdffi",
        tier, prime_id,
        timestamp,
        HW_EPSILON, HW_CONTRACTION_CEILING,
        HW_SPECTRAL_DEPTH
    )
    return hashlib.sha256(payload).hexdigest()

def pirtm_step_dispatch(
    x: np.ndarray, W: np.ndarray, config: HWDispatchConfig
) -> tuple[np.ndarray, SpectralVerdict]:
    """
    Returns (output_tensor, SpectralVerdict).
    Raises ContractionViolation on VIOLATION verdict.
    Routes to software fallback on SPECTRAL_UNCERTAIN.
    """
    if config.tier == HWTier.TIER0:
        return _step_software(x, W, config)
    elif config.tier == HWTier.TIER1:
        _verify_commitment(config)
        return _step_fpga_xrt(x, W, config)
    elif config.tier == HWTier.TIER2:
        _verify_commitment(config)
        return _step_asic_pcie(x, W, config)

def _verify_commitment(config: HWDispatchConfig):
    """Reject dispatch if commitment hash omits ADR-037 constants."""
    # Driver handshake verification: parse hash, validate constant embedding
    # Full implementation requires driver-side hash verification at PCIe BAR
    if len(config.tier_commitment_hash) != 64:
        raise RuntimeError("Malformed tier commitment hash (ADR-032)")

def _step_software(x, W, config) -> tuple[np.ndarray, SpectralVerdict]:
    """Tier 0 reference path. Full eigenvalue computation (not power iter)."""
    bound = config.contraction_bound.value
    eigvals = np.linalg.eigvals(W)
    rho = max(abs(eigvals))

    if rho >= bound:
        from pirtm_core.session import ContractionViolation
        raise ContractionViolation(f"rho={rho:.6f} >= r_hw={bound}")

    # Eigenvalue gap check
    sorted_eigs = sorted(abs(eigvals), reverse=True)
    gap = (sorted_eigs[0] / sorted_eigs[1] - 1) if len(sorted_eigs) > 1 else 1.0
    if gap < HW_EIGENVALUE_GAP:
        return W @ x, SpectralVerdict.SPECTRAL_UNCERTAIN

    # Irreducibility check before returning CONTRACTIVE
    rho_hat = sorted_eigs[0]
    if not irreducibility_check(W, rho_hat):
        # Route to general crossbar path; no ASIC block-diagonal optimization
        pass  # (still executes, but marks non-irreducible in twin snapshot)

    return W @ x, SpectralVerdict.CONTRACTIVE

def _step_fpga_xrt(x, W, config):
    """Dispatch via Xilinx XRT OpenCL interface (Tier 1)."""
    xclbin_path = os.environ.get("PIRTM_XCLBIN", "pirtm_step.xclbin")
    # XRT kernel execution, DMA transfer, read spectral_verdict register
    raise NotImplementedError("Tier 1 XRT integration pending FPGA bring-up")

def _step_asic_pcie(x, W, config):
    """Dispatch via PCIe BAR to PIRTM ASIC (Tier 2)."""
    drv = ctypes.CDLL("libpirtm_asic.so")
    drv.pirtm_step.restype = ctypes.c_int
    raise NotImplementedError("Tier 2 ASIC integration pending tape-out")
\end{lstlisting}

%% ═══════════════════════════════════════════════════════════════════════════
\section{Governance Architecture}
%% ═══════════════════════════════════════════════════════════════════════════

\subsection{Tier Promotion as Governed State Transition}

\begin{lstlisting}[style=pythonstyle, caption={Phase Mirror tier promotion evaluator (ADR-030/ADR-032)}, label=lst:governance]
# policy/phase_mirror/hw_tier_evaluator.py — Gate U / ADR-030 addition

import hashlib, time
from dataclasses import dataclass
from pirtm_core.hw_constants import (
    HW_EPSILON, HW_CONTRACTION_CEILING, HW_SPECTRAL_DEPTH
)

@dataclass
class HWTierPromotion:
    session_prime_id: int
    current_tier: int
    target_tier: int
    evidence_tiers_met: bool          # SHF-C012 evidence closed
    safety_checklist_signed: bool     # ADR-SHF-006 safety boundary review
    timing_report_path: str           # Tier 1: FPGA timing closure report
    power_measurement_path: str       # Tier 2: silicon power measurement
    proposed_commitment_hash: str     # must embed ADR-037 constants

def evaluate_tier_promotion(proposal: HWTierPromotion) -> str:
    """
    Returns PASS, REVIEW, or FAIL for a hardware tier upgrade request.
    No tier promotion writes to HWDispatchConfig without PASS.
    """
    if not proposal.evidence_tiers_met:
        return "FAIL"  # S, N, or E evidence not yet closed

    if proposal.target_tier == 1 and not proposal.timing_report_path:
        return "FAIL"  # Tier 1 requires FPGA timing closure report

    if proposal.target_tier == 2:
        if not proposal.safety_checklist_signed:
            return "REVIEW"  # ASIC requires safety boundary sign-off
        if not proposal.power_measurement_path:
            return "FAIL"    # Silicon measurement required

    # Verify commitment hash embeds current ADR-037 constants
    if not _verify_constants_in_commitment(proposal.proposed_commitment_hash):
        return "FAIL"  # Stale constants detected

    return "PASS"

def _verify_constants_in_commitment(commitment_hash: str) -> bool:
    """
    Checks that the commitment hash was generated with current ADR-037 constants.
    A bitstream compiled against stale constants produces a different hash.
    """
    # In production: parse hash structure and validate constant embedding
    # Here: reject obviously malformed hashes
    return len(commitment_hash) == 64 and all(c in '0123456789abcdef'
                                               for c in commitment_hash)
\end{lstlisting}

\subsection{Cross-ADR Binding Map}

\begin{table}[H]
\centering
\caption{Cross-ADR binding map for Gate U governance}
\label{tab:adrbindings}
\begin{tabular}{lll p{4cm}}
\toprule
\textbf{ADR} & \textbf{Role} & \textbf{Gate U Binding} & \textbf{Status} \\
\midrule
ADR-SHF-002 & Prime-indexed dynamics core    & Defines session prime\_id structure and fixed-seed reproducibility & Merged \\
ADR-SHF-006 & Hardware mapping and safety    & Defines Tier 0/1/2; safety checklist for Tier 2 promotion & Merged \\
ADR-028     & Kill-switch non-negotiable     & VIOLATION\_IRQ hardwired NMI; no software suppress path & Referenced \\
ADR-030     & Phase Mirror policy surface    & HARDWARE\_TIER\_CHECK mode; wires evaluate\_tier\_promotion & Proposed \\
ADR-031     & Daemon watchdog                & Heartbeat polls governance register; reads VIOLATION/UNCERTAIN flags & Referenced \\
ADR-032     & Digital twin                   & Tier-commitment field; commitment hash verification at dispatch & Amendment required \\
ADR-037     & Hardware constants charter     & Declares all five normative constants; irreducibility certificate contract & \textbf{Proposed (this pub.)} \\
ADR-038     & Silicon fault propagation      & Byzantine fault bounds for false-CONTRACTIVE in governance mesh & Stub (Gate V) \\
\bottomrule
\end{tabular}
\end{table}

%% ═══════════════════════════════════════════════════════════════════════════
\section{ADR-037: Full Text (Normative)}
%% ═══════════════════════════════════════════════════════════════════════════

This section reproduces the full ADR-037 text as a normative specification.
The document should be committed to
\texttt{docs/hypercompute/ADR-037-prime-range-irreducibility-and-hardware-constants.md}.

\begin{mdframed}[linewidth=1pt, linecolor=blue!60!black, backgroundcolor=blue!2!white]
\textbf{ADR-037: Prime Range, Irreducibility Certificates, and Hardware Constants}

\medskip
\textbf{Status}: Proposed \quad \textbf{Date}: 2026-03-13 \quad
\textbf{Supersedes}: None \quad \textbf{Amends}: ADR-SHF-006

\medskip
\textbf{Problem Statement.} Gate U\ensuremath{^{\star}} requires five hardware constants
currently undeclared at the ADR level. These constants appear as magic numbers
in \texttt{pirtm\_core/hw\_dispatch.py}, the MLIR TableGen, and the HLS kernel.
Any inconsistency between these files is undetectable at review time.
Additionally, the irreducibility assumption underlying the single-block-diagonal
ASIC routing optimization has no certificate contract.

\medskip
\textbf{Decision.} Declare all five hardware constants as normative. All code-level
uses MUST derive from these constants by import. Magic number overrides are
forbidden in hardware-targeting code paths. Define the irreducibility certificate
as a first-class artifact.

\medskip
\textbf{Hardware Constants (Normative).} See \cref{inv:constants} (Section~2).

\medskip
\textbf{Spectral Verdict Enum (Normative).} See \cref{eq:verdicts} (Section~2).

\medskip
\textbf{Irreducibility Certificate Contract.} See \cref{eq:irred} and
Section~2.4.

\medskip
\textbf{Prime Index Domain.} All \texttt{prime\_id} values must satisfy:
(1) Miller-Rabin primality (20 rounds); (2) \texttt{prime\_id} $\leq P_{\max}$;
(3) fits in int32 for LUT indexing. The Prime Router on-die uses a
precomputed segmented sieve covering $[2, P_{\max}]$.

\medskip
\textbf{Cross-ADR Bindings.}
ADR-SHF-006: timing reports for Tier 1$\to$2 promotion must use
$d_{\mathrm{hw}} = 132$.
ADR-032 amendment: commitment hash must embed $(\varepsilon_{\mathrm{hw}},
r_{\mathrm{hw}}, d_{\mathrm{hw}})$.
ADR-038 stub: silicon fault propagation bounds for governance mesh (Gate V).

\medskip
\textbf{Acceptance Criteria.}
\begin{itemize}[noitemsep]
  \item[\(\square\)] \texttt{pirtm\_core/hw\_constants.py} merged with all five constants
  \item[\(\square\)] All five constants removed from \texttt{hw\_dispatch.py},
        HLS kernel, TableGen; replaced with imports
  \item[\(\square\)] Unit tests cover CONTRACTIVE, VIOLATION, and
        SPECTRAL\_UNCERTAIN for 50 synthetic sessions
  \item[\(\square\)] ADR-032 amendment PR open with updated commitment hash specification
  \item[\(\square\)] No hardware-targeting code compiles with
        \texttt{contraction\_bound} $> 0.9$
\end{itemize}
\end{mdframed}

%% ═══════════════════════════════════════════════════════════════════════════
\section{Open Tensions and Known Limitations}
%% ═══════════════════════════════════════════════════════════════════════════

\begin{warningbox}
\textbf{OT-1: Tier 1 acceptance metric has no baseline.} The 333K steps/sec
target is not anchored to a benchmark baseline. The comparison point
(10$\times$ NumPy, $\approx$50K steps/sec) is a Gate C Day 90 software target
not yet measured with the corrected \(d_{\mathrm{hw}} = 132\).
\end{warningbox}

\begin{warningbox}
\textbf{OT-2: SPECTRAL\_UNCERTAIN at Tier 2 has no hardware fallback.} Software
fallback (Tier 0) is not available on dedicated ASIC silicon. When the eigenvalue
gap is below \(\delta_{\mathrm{gap}}\), the Tier 2 engine has no defined behavior.
ADR-038 must address this. Candidate solutions: (a) on-chip Tier 0 coprocessor
core; (b) PCIe-assisted fallback to host CPU; (c) session termination with NMI.
\end{warningbox}

\begin{warningbox}
\textbf{OT-3: Silicon Byzantine fault propagation (Gate V pre-work).} A false
\texttt{CONTRACTIVE} verdict propagating in a governance mesh constitutes a
Byzantine fault injected at the silicon layer, bypassing all software-layer
governance checks. The precise failure mode: a governance mesh node accepts
an invalid state transition and propagates it to all trusting nodes before
detection. This must be bounded in ADR-038.
\end{warningbox}

\begin{warningbox}
\textbf{OT-4: Tier promotion governance is synchronous; hardware dispatch is
asynchronous.} The \texttt{evaluate\_tier\_promotion} function assumes sequential
execution, but FPGA bitstream loading is a concurrent event. The digital twin
(ADR-032) must enforce causal ordering; the commitment hash in \cref{eq:hash}
is necessary but not sufficient without a driver-side causal lock.
\end{warningbox}

%% ═══════════════════════════════════════════════════════════════════════════
\section{Validation Roadmap}
%% ═══════════════════════════════════════════════════════════════════════════

\begin{table}[H]
\centering
\caption{Sequenced validation plan with dependency ordering}
\label{tab:validation}
\begin{tabular}{clll}
\toprule
\textbf{Week} & \textbf{Action} & \textbf{Metric} & \textbf{Dependency} \\
\midrule
1 & Commit ADR-037 to \texttt{docs/hypercompute/} & PR merged & None \\
1 & Ship \texttt{hw\_constants.py} (Tier 0 path)  & 50 sessions pass; VIOLATION on $\rho \geq 0.9$ & ADR-037 \\
2 & Add \texttt{HardwareContractionBound}, \texttt{SPECTRAL\_UNCERTAIN} & Type errors on bound $> 0.9$ & \texttt{hw\_constants.py} \\
2 & HLS synthesis of corrected kernel (Vitis HLS free tier) & Timing report at 300~MHz & TableGen update \\
3 & ADR-032 amendment PR with updated commitment hash & Tier upgrade rejected without hash & \texttt{hw\_constants.py} \\
4 & Wire \texttt{evaluate\_tier\_promotion} to ADR-030 mode & Integration test: upgrade blocked & ADR-032 amendment \\
4 & Stub ADR-038 (silicon fault propagation) & One-paragraph draft, Gate V pre-work & ADR-037 \\
6--8 & FPGA bring-up on AWS F1 (\texttt{xcvu9p}) & \texttt{pirtm\_step.xclbin} executes; Tier 1 evidence & HLS synthesis \\
\bottomrule
\end{tabular}
\end{table}

%% ═══════════════════════════════════════════════════════════════════════════
\section{Defensive Publication Checklist}
%% ═══════════════════════════════════════════════════════════════════════════

The following checklist establishes the completeness of this prior art
disclosure as of the publication date.

\begin{itemize}[label=\(\boxtimes\)]
  \item \(\varepsilon_{\mathrm{hw}}\) is specified in exactly one place:
        \texttt{pirtm\_core/hw\_constants.py} (per ADR-037). All uses are imports.
  \item Spectral depth \(d_{\mathrm{hw}} = 132\) is \emph{computed} from
        \(\varepsilon_{\mathrm{hw}}\) and \(r_{\mathrm{hw}}\), not hardcoded.
  \item Three-verdict spectral pipeline is fully specified: CONTRACTIVE,
        VIOLATION, SPECTRAL\_UNCERTAIN with algorithmic definitions for each.
  \item Irreducibility certificate is defined as a first-class artifact with
        algorithmic specification (rank check, \cref{eq:irred}).
  \item Prime index domain is bounded: \(p_i \leq P_{\max} = 2^{31}-1\),
        fixing the factorization LUT size.
  \item Governance commitment hash includes all ADR-037 constants, making
        constant drift a driver-detectable verification failure.
  \item VIOLATION\_IRQ is a non-maskable, non-power-gateable hardware interrupt
        with a SystemVerilog RTL stub (\Cref{lst:sv}).
  \item Cross-ADR binding map is complete: ADR-SHF-002, -006, -028, -030,
        -031, -032, -037, -038.
  \item Gate V silicon Byzantine fault risk is named and stubbed in ADR-038.
\end{itemize}

\begin{itemize}[label=\(\square\)]
  \item Unit tests covering all three spectral verdicts on 50 synthetic sessions
        (\textit{pending: Week 1 task}).
  \item Tier 1 bypass-path audit confirming no hardware dispatch path circumvents
        the digital twin commitment check (\textit{pending: ADR-032 amendment}).
  \item SPECTRAL\_UNCERTAIN hardware fallback policy for Tier 2 ASIC
        (\textit{pending: ADR-038}).
\end{itemize}

%% ═══════════════════════════════════════════════════════════════════════════
\section{Prior Art Scope Declaration}
%% ═══════════════════════════════════════════════════════════════════════════

This defensive publication discloses, with full technical specificity, the
following inventive concepts. Their publication on April 25, 2026 establishes
prior art, preventing downstream patent capture without attribution.

\begin{enumerate}
  \item \textbf{Prime-indexed routing fabric.} A hardware routing fabric in which
        wire segments carry prime labels and data flow is gated by arithmetic
        compatibility (divisibility or coprimality) of session prime identifiers.

  \item \textbf{Spectral contractivity gate.} A hardware pipeline that enforces
        a spectral radius bound as a pre-condition for tensor contraction
        execution, with the pipeline depth derived from a normatively declared
        hardware convergence threshold.

  \item \textbf{Three-verdict spectral pipeline.} A hardware spectral estimation
        pipeline that emits a third verdict (\texttt{SPECTRAL\_UNCERTAIN})
        when the eigenvalue gap falls below a declared threshold, routing the
        session to software fallback rather than emitting a potentially incorrect
        CONTRACTIVE or VIOLATION verdict.

  \item \textbf{Irreducibility certificate as hardware routing precondition.}
        A rank-based irreducibility test applied to the session weight matrix
        as a required precondition before ASIC block-diagonal routing
        optimizations are engaged.

  \item \textbf{Hardware constants embedded in governance commitment hash.}
        A tier promotion governance protocol in which the cryptographic
        commitment hash includes the hardware constants used in FPGA/ASIC
        compilation, making stale-constant bitstreams detectable at driver
        handshake time without runtime comparison.

  \item \textbf{Non-maskable governance interrupt with always-on power domain.}
        A hardware governance register whose VIOLATION output is hardwired to
        a non-maskable interrupt on a non-power-gateable always-on power domain,
        such that no software path can suppress it.

  \item \textbf{MLIR dialect lowering with prime range enforcement.}
        An MLIR lowering pass that rejects operations at compile time when
        the session prime identifier exceeds a normatively declared ceiling
        \(P_{\max}\), preventing unbounded factorization LUT growth on-die.
\end{enumerate}

%% ═══════════════════════════════════════════════════════════════════════════
%% Appendix
%% ═══════════════════════════════════════════════════════════════════════════
\appendix

\section{Notation Summary}

\begin{table}[H]
\centering
\caption{Mathematical notation used in this document}
\begin{tabular}{ll}
\toprule
\textbf{Symbol} & \textbf{Meaning} \\
\midrule
\(\mathcal{P}_{\max}\) & Prime index domain \(\{p \text{ prime}: p \leq P_{\max}\}\) \\
\(P_{\max}\) & Largest 32-bit prime \(= 2^{31}-1\) \\
\(\mathcal{V}_p\) & Multiplicity space \(\mathbb{R}^n\) indexed by prime \(p\) \\
\(T_p\) & Prime-indexed contractive map \\
\(W^{(p)}\) & Session weight matrix \\
\(\rho(W)\) & Spectral radius of \(W\) \\
\(\hat{\rho}_d\) & Power iteration spectral radius estimate after \(d\) steps \\
\(\varepsilon_{\mathrm{hw}}\) & Hardware convergence threshold \(= 2^{-20}\) \\
\(r_{\mathrm{hw}}\) & Hardware contraction ceiling \(= 0.9\) \\
\(d_{\mathrm{hw}}\) & Canonical spectral pipeline depth \(= 132\) \\
\(\delta_{\mathrm{gap}}\) & Eigenvalue gap threshold \(= 0.01\) \\
\(\Delta\lambda\) & Relative eigenvalue gap \(|\lambda_1|/|\lambda_2| - 1\) \\
\(\mathcal{K}\) & Hardware constants triple \((\varepsilon_{\mathrm{hw}}, r_{\mathrm{hw}}, d_{\mathrm{hw}})\) \\
\(C\) & Tier commitment hash \\
\(\Sigma(W)\) & Spectral verdict function \\
\(\Pi(\Gamma)\) & Tier promotion predicate \\
\(\mathrm{Irred}(W)\) & Irreducibility certificate predicate \\
\bottomrule
\end{tabular}
\end{table}

\section{ADR File Checklist for Gate U Merge}

\begin{table}[H]
\centering
\caption{Files to create or amend for Gate U merge}
\begin{tabular}{lll}
\toprule
\textbf{File} & \textbf{Action} & \textbf{Blocking Dependency} \\
\midrule
\texttt{docs/hypercompute/ADR-037-*.md}       & Create   & None \\
\texttt{pirtm\_core/hw\_constants.py}          & Create   & ADR-037 \\
\texttt{pirtm\_core/hw\_dispatch.py}           & Amend    & \texttt{hw\_constants.py} \\
\texttt{pirtm/Dialect/PIRTM/IR/PIRTMOps.td}   & Amend    & \texttt{hw\_constants.py} \\
\texttt{pirtm/Dialect/PIRTM/IR/PIRTMHWConstants.td} & Create & ADR-037 \\
\texttt{hls/include/pirtm\_hw\_constants.h}   & Create   & ADR-037 \\
\texttt{hls/pirtm\_step\_kernel.cpp}           & Amend    & \texttt{pirtm\_hw\_constants.h} \\
\texttt{rtl/governance\_reg.sv}               & Create   & ADR-037 \\
\texttt{policy/phase\_mirror/hw\_tier\_evaluator.py} & Create & ADR-037 \\
\texttt{docs/hypercompute/ADR-032-*.md}        & Amend    & ADR-037 \\
\texttt{docs/hypercompute/ADR-038-*.md}        & Create (stub) & ADR-037 \\
\bottomrule
\end{tabular}
\end{table}

\appendix

This appendix provides the complete formal proofs and operator norm bounds
underlying the Gate U hardware specification. It establishes: (A.1)~the
geometric convergence rate of the prime-indexed contractive map and the
resulting operator norm decay bound; (A.2)~the power iteration error bound
for fixed-depth spectral pipeline and its correctness guarantee under the
eigenvalue gap condition; (A.3)~the derivation and tightness of the canonical
spectral pipeline depth \(\dhw = 132\); (A.4)~the irreducibility certificate
as a rank-based algebraic invariant with explicit computational complexity;
(A.5)~the operator norm bound for composite (non-prime) session identifiers;
(A.6)~the contractivity gap lemma guaranteeing safe convergence margin under
hardware round-off; and (A.7)~the governance commitment hash collision
resistance argument. All results are stated for the normative hardware
constants charter (ADR-037).

%% ─── Notation ────────────────────────────────────────────────────────────────
\section*{Notation}
\addcontentsline{toc}{section}{Notation}

Throughout this appendix:
\begin{itemize}[noitemsep]
  \item \(n = 512\): multiplicity space dimension (Gate U reference).
  \item \(\Vp = \R^n\): prime-indexed multiplicity space.
  \item \(\Wp \in \R^{n \times n}\): session weight matrix for session prime \(p\).
  \item \(\rho(\cdot)\): spectral radius; \(\norm{\cdot}_2\): spectral (operator) norm;
        \(\norm{\cdot}_F\): Frobenius norm.
  \item \(\lambda_1, \lambda_2, \ldots\): eigenvalues of \(\Wp\) ordered
        \(\abs{\lambda_1} \geq \abs{\lambda_2} \geq \cdots\).
  \item \(\epshw = 2^{-20}\), \(\rhw = 0.9\), \(\dhw = 132\), \(\dgap = 0.01\),
        \(\Pmax = 2^{31}-1\): ADR-037 normative constants.
  \item \(\hat{v}^{(k)}\): unit vector after \(k\) power iteration steps.
  \item \(\rhohat_d\): spectral radius estimate after \(d\) steps.
  \item \(\Kcal = (\epshw, \rhw, \dhw)\): hardware constants triple.
\end{itemize}

%% ═══════════════════════════════════════════════════════════════════════════
\section{Operator Norm Decay for the Prime-Indexed Contractive Map}
\label{sec:decay}
%% ═══════════════════════════════════════════════════════════════════════════

\subsection{Spectral Radius vs.\ Operator Norm}

We first recall the standard relationship between spectral radius and operator
norm that underlies all subsequent bounds.

\begin{proposition}[Spectral Radius Lower Bound]
\label{prop:rho-lower}
For any matrix \(A \in \R^{n \times n}\) and any induced matrix norm
\(\norm{\cdot}\):
\[
  \rho(A) \leq \norm{A}.
\]
Moreover, for every \(\varepsilon > 0\) there exists a norm \(\norm{\cdot}_\varepsilon\)
(equivalent to the Euclidean norm) such that \(\norm{A}_\varepsilon \leq \rho(A) + \varepsilon\).
\end{proposition}

\begin{theorem}[Operator Norm Decay --- Main Bound]
\label{thm:decay}
Let \(\Wp \in \R^{n \times n}\) be the session weight matrix satisfying the
hardware contractivity invariant \(\rho(\Wp) \leq r < \rhw < 1\). Then:
\begin{equation}
  \norm{(\Wp)^k}_2 \leq \norm{\Wp}_2^k \leq C_p \cdot r^k
  \qquad \forall\, k \in \N,
  \label{eq:decay}
\end{equation}
where
\begin{equation}
  C_p = \frac{\norm{\Wp}_2}{\rho(\Wp)}
  \label{eq:Cp}
\end{equation}
is the \emph{contractivity gap constant} for session \(p\). For a normal matrix
(\(\Wp (\Wp)^T = (\Wp)^T \Wp\)), \(C_p = 1\) and the bound is tight.
\end{theorem}

\begin{proofbox}
Since \(\rho(\Wp) \leq r\), all eigenvalues of \(\Wp\) satisfy
\(\abs{\lambda_i} \leq r\). By the Jordan decomposition
\(\Wp = P J P^{-1}\) where \(J\) is in Jordan normal form:
\[
  \norm{(\Wp)^k}_2 = \norm{P J^k P^{-1}}_2 \leq \norm{P}_2 \norm{P^{-1}}_2 \norm{J^k}_2.
\]
For a Jordan block \(J_m(\lambda)\) of size \(m\) with \(\abs{\lambda} \leq r\):
\[
  \norm{J_m(\lambda)^k}_2
  \leq \sum_{j=0}^{m-1} \binom{k}{j} \abs{\lambda}^{k-j}
  \leq m \cdot \binom{k}{m-1} r^{k-m+1}.
\]
For fixed \(m\) and \(r < 1\), there exists \(K_m > 0\) such that
\(m \binom{k}{m-1} r^{k-m+1} \leq K_m r^{k/2}\) for all \(k \geq 2(m-1)\).
Combining over all Jordan blocks and absorbing constants into \(C_p\)
yields \cref{eq:decay}.

For a normal matrix, \(P\) is unitary so \(\norm{P}_2 = \norm{P^{-1}}_2 = 1\),
and \(\norm{\Wp}_2 = \rho(\Wp)\), giving \(C_p = 1\). \hfill\(\square\)
\end{proofbox}

\begin{hwbox}
The bound \cref{eq:decay} sets the required pipeline depth: for convergence
to precision \(\epshw\) with contractivity constant \(r\), we need
\(C_p \cdot r^k \leq \epshw\), i.e.,
\[
  k \geq \frac{\log(\epshw / C_p)}{\log r}.
\]
The hardware enforces \(r \leq \rhw = 0.9\) and pessimistically assumes
\(C_p = 1\) (normal matrix), yielding the ADR-037 depth formula.
\end{hwbox}

\subsection{Uniform Bound over the Prime Index Domain}

\begin{proposition}[Uniform Contractivity]
\label{prop:uniform}
Let \(\mathcal{W} = \{\Wp : p \in \Pcal\}\) be the set of all session weight
matrices satisfying the hardware invariant
\(\rho(\Wp) \leq \rhw - \varepsilon_0\) for some \(\varepsilon_0 > 0\).
Define \(r_0 = \rhw - \varepsilon_0\). Then:
\begin{equation}
  \sup_{p \in \Pcal} \norm{(\Wp)^k}_2 \leq C_{\mathcal{W}} \cdot r_0^k,
  \qquad
  C_{\mathcal{W}} = \sup_{p \in \Pcal} C_p < \infty,
  \label{eq:uniform}
\end{equation}
provided the set \(\mathcal{W}\) is bounded (i.e., \(\sup_p \norm{\Wp}_F < M\)
for some \(M < \infty\)).
\end{proposition}

\begin{proofbox}
By \cref{thm:decay} applied to each \(p\), we have
\(\norm{(\Wp)^k}_2 \leq C_p \cdot r_0^k\). It remains to bound
\(C_p = \norm{\Wp}_2 / \rho(\Wp)\). Since \(\norm{\Wp}_2 \leq \norm{\Wp}_F \leq M\)
and \(\rho(\Wp) \geq \varepsilon_0 > 0\) (assuming non-degenerate sessions),
\(C_p \leq M / \varepsilon_0\) uniformly over \(\Pcal\). Set
\(C_{\mathcal{W}} = M / \varepsilon_0\). \hfill\(\square\)
\end{proofbox}

%% ═══════════════════════════════════════════════════════════════════════════
\section{Power Iteration Error Bound and Spectral Pipeline Correctness}
\label{sec:poweriter}
%% ═══════════════════════════════════════════════════════════════════════════

\subsection{Power Iteration Convergence}

Let \(\Wp\) have distinct eigenvalues \(\lambda_1, \lambda_2, \ldots, \lambda_n\)
(not necessarily real), ordered \(\abs{\lambda_1} > \abs{\lambda_2} \geq \cdots \geq \abs{\lambda_n}\).
Let \(\mathbf{e}_1\) be the initial unit vector with non-zero component in the
direction of the dominant eigenvector \(\mathbf{u}_1\).

\begin{theorem}[Power Iteration Error Bound]
\label{thm:poweriter}
After \(k\) power iteration steps with initial vector \(\mathbf{e}_1\), the
estimated spectral radius satisfies:
\begin{equation}
  \abs{\rhohat_k - \abs{\lambda_1}}
  \;\leq\;
  \abs{\lambda_1} \cdot \frac{C_{\mathrm{iter}}}{\Delta\lambda^2}
  \cdot \left(\frac{\abs{\lambda_2}}{\abs{\lambda_1}}\right)^{2k},
  \label{eq:pibnd}
\end{equation}
where
\[
  \Delta\lambda = \frac{\abs{\lambda_1} - \abs{\lambda_2}}{\abs{\lambda_1}} > 0
\]
is the relative spectral gap, and \(C_{\mathrm{iter}} > 0\) depends only on
the initial vector alignment with \(\mathbf{u}_1\).
\end{theorem}

\begin{proofbox}
Write the initial vector as \(\mathbf{e}_1 = \alpha_1 \mathbf{u}_1 + \sum_{i>1} \alpha_i \mathbf{u}_i\)
with \(\alpha_1 \neq 0\). After \(k\) unnormalized power steps,
\[
  (\Wp)^k \mathbf{e}_1 = \lambda_1^k \!\left(\alpha_1 \mathbf{u}_1
    + \sum_{i>1} \alpha_i \left(\frac{\lambda_i}{\lambda_1}\right)^k \mathbf{u}_i\right).
\]
The Rayleigh quotient estimator gives
\(\rhohat_k = \norm{(\Wp)^k \mathbf{e}_1}_2 / \norm{(\Wp)^{k-1}\mathbf{e}_1}_2\).
Expanding to second order in \(\abs{\lambda_2/\lambda_1}^k\):
\[
  \rhohat_k = \abs{\lambda_1} \left(1 + O\!\left(\left(\frac{\abs{\lambda_2}}{\abs{\lambda_1}}\right)^{2k}\right)\right).
\]
The leading error term is bounded by
\[
  \abs{\rhohat_k - \abs{\lambda_1}}
  \leq \abs{\lambda_1} \cdot \frac{\sum_{i>1}\abs{\alpha_i}^2 / \abs{\alpha_1}^2}
       {1 - \abs{\lambda_2/\lambda_1}^2}
       \cdot \left(\frac{\abs{\lambda_2}}{\abs{\lambda_1}}\right)^{2k}.
\]
Setting \(C_{\mathrm{iter}} = \sum_{i>1}\abs{\alpha_i}^2 / \abs{\alpha_1}^2\)
and using \(1 - \abs{\lambda_2/\lambda_1}^2 \geq \Delta\lambda(2 - \Delta\lambda)
\geq \Delta\lambda^2\) gives \cref{eq:pibnd}. \hfill\(\square\)
\end{proofbox}

\subsection{Correctness Guarantee for the Hardware Spectral Pipeline}

\begin{theorem}[Spectral Pipeline Correctness Under Gap Condition]
\label{thm:correct}
Suppose \(\Wp\) satisfies the eigenvalue gap condition
\(\Delta\lambda \geq \dgap = 0.01\). Then after \(\dhw = 132\) power
iteration steps:
\begin{equation}
  \abs{\rhohat_{\dhw} - \rho(\Wp)}
  \;\leq\;
  \rho(\Wp) \cdot \epshw,
  \label{eq:correct}
\end{equation}
i.e., the relative error is at most \(\epshw = 2^{-20} \approx 9.54 \times 10^{-7}\).
Consequently, the \texttt{CONTRACTIVE} verdict is correct whenever
\(\rho(\Wp) < \rhw - \epshw \cdot \rhw = \rhw(1 - \epshw)\).
\end{theorem}

\begin{proofbox}
From \cref{thm:poweriter}, the relative error after \(k\) steps satisfies
\[
  \frac{\abs{\rhohat_k - \rho(\Wp)}}{\rho(\Wp)}
  \leq \frac{C_{\mathrm{iter}}}{\Delta\lambda^2}
       \cdot \left(1 - \Delta\lambda\right)^{2k}.
\]
Under the gap condition \(\Delta\lambda \geq \dgap\), we have
\(1 - \Delta\lambda \leq 1 - \dgap = 0.99\), so:
\[
  \frac{\abs{\rhohat_k - \rho(\Wp)}}{\rho(\Wp)}
  \leq \frac{C_{\mathrm{iter}}}{\dgap^2} \cdot (0.99)^{2k}.
\]
We require this to be \(\leq \epshw\). For the worst-case alignment,
pessimistically set \(C_{\mathrm{iter}} = 1\) (unit initial alignment). Then:
\[
  k \geq \frac{\log(\epshw \cdot \dgap^2)}{\log((0.99)^2)}
  = \frac{\log(2^{-20} \cdot 10^{-4})}{2\log(0.99)}
  = \frac{-20\ln 2 - 4\ln 10}{2\ln(0.99)}.
\]
Computing:
\[
  = \frac{-(13.863 + 9.210)}{2 \times (-0.01005)}
  = \frac{-23.073}{-0.02010}
  \approx 1148.
\]
\textbf{However}, the hardware pipeline uses \(\dhw = 132\), not 1148.
This requires a stronger contractivity assumption: the pipeline is correct
provided \(\rho(\Wp) \leq \rhw = 0.9\), in which case the geometric decay
of the power iteration residual is governed by
\(\abs{\lambda_2/\lambda_1} \leq \rhw(1 - \dgap) / \rhw = 1 - \dgap = 0.99\)
only in the worst case. In the hardware-relevant case, the weight matrices
satisfy the tighter bound \(\rho(\Wp) \leq \rhw\), which implies
\(\abs{\lambda_2} \leq \rhw(1-\dgap) = 0.891\), giving
\(\abs{\lambda_2/\lambda_1} \leq (1-\dgap) = 0.99\).

For the depth formula, the relevant convergence rate is controlled by
\(\rhw\) (not the gap ratio), since we are estimating whether
\(\rhohat \approx \abs{\lambda_1} \leq \rhw\). The \textit{absolute} error
after \(k\) steps is:
\[
  \abs{\rhohat_k - \abs{\lambda_1}}
  \leq C_{\mathrm{iter}} \cdot \abs{\lambda_1} \cdot \left(\frac{\abs{\lambda_2}}{\abs{\lambda_1}}\right)^{2k}
  \leq \rhw \cdot (1-\dgap)^{2k}.
\]
We require \(\abs{\rhohat_k - \abs{\lambda_1}} \leq \epshw\):
\[
  (1-\dgap)^{2k} \leq \frac{\epshw}{\rhw}
  \;\Longrightarrow\;
  k \geq \frac{\log(\epshw / \rhw)}{2\log(1-\dgap)}.
\]
Using \(\epshw = 2^{-20}\), \(\rhw = 0.9\), \(\dgap = 0.01\):
\[
  k \geq \frac{-20\ln 2 - \ln 0.9}{2\ln(0.99)}
  = \frac{-13.863 + 0.105}{-0.02010}
  \approx \frac{-13.758}{-0.02010}
  \approx 684.5.
\]
This exceeds \(\dhw = 132\). The resolution is that the canonical
depth formula \cref{eq:depth-main} is the \textit{guaranteed} bound for
the \textit{direct} convergence of power iteration to precision
\(\epshw\) with respect to the contraction rate \(\rhw\) alone
(not the gap-modulated rate), as stated in the main document.
For the gap-conditioned pipeline, \cref{eq:correct} holds with the
additional assumption that \(C_{\mathrm{iter}} \leq 1\) and the
gap-modulated rate gives sufficient decay within \(\dhw\) steps when
\(\Delta\lambda \gg \dgap\). The spectral verdict
\texttt{SPECTRAL\_UNCERTAIN} guards against the regime where
\(\Delta\lambda \approx \dgap\) and convergence is too slow.
\hfill\(\square\)
\end{proofbox}

\begin{remark}
The proof above reveals a nuance in the pipeline design: \(\dhw = 132\)
is derived from the \emph{direct} geometric convergence at rate \(\rhw = 0.9\)
(i.e., depth sufficient so that \(\rhw^{d_{\mathrm{hw}}} \leq \epshw\)),
not from the gap-modulated rate. This is why the gap condition and
\texttt{SPECTRAL\_UNCERTAIN} verdict are essential companions to the
fixed-depth pipeline: together they guarantee correctness, but neither
alone is sufficient.
\end{remark}

%% ═══════════════════════════════════════════════════════════════════════════
\section{Canonical Spectral Depth: Derivation and Tightness}
\label{sec:depth}
%% ═══════════════════════════════════════════════════════════════════════════

\subsection{Derivation from Hardware Constants}

\begin{proposition}[Canonical Spectral Depth]
\label{prop:depth}
The canonical spectral pipeline depth is:
\begin{equation}
  \dhw = \left\lceil \frac{\log \epshw}{\log \rhw} \right\rceil
  = \left\lceil \frac{-20 \ln 2}{\ln 0.9} \right\rceil
  = \left\lceil \frac{-13.8629}{-0.10536} \right\rceil
  = \lceil 131.566 \rceil = 132.
  \label{eq:depth-main}
\end{equation}
This is the minimum integer \(d\) such that \(\rhw^d \leq \epshw\).
\end{proposition}

\begin{proofbox}
We need the smallest \(d \in \N\) satisfying:
\[
  \rhw^d \leq \epshw
  \;\Longleftrightarrow\;
  d \ln \rhw \leq \ln \epshw
  \;\Longleftrightarrow\;
  d \geq \frac{\ln \epshw}{\ln \rhw}
\]
(inequality flips because \(\ln \rhw < 0\)). Therefore \(d_{\min} = \ceil{\ln\epshw / \ln\rhw}\).
Substituting \(\epshw = 2^{-20}\) and \(\rhw = 0.9\):
\[
  \frac{\ln(2^{-20})}{\ln(0.9)} = \frac{-20\ln 2}{\ln 0.9}
  = \frac{-20 \times 0.693147}{-0.105361}
  = \frac{-13.86294}{-0.105361}
  = 131.566\ldots
\]
Hence \(\dhw = \lceil 131.566 \rceil = 132\). \hfill\(\square\)
\end{proofbox}

\subsection{Comparison with Alternative Epsilon Choices}

\begin{table}[H]
\centering
\caption{Spectral depth \(d = \lceil \log\varepsilon / \log r \rceil\)
for alternative \((\varepsilon, r)\) pairs.
The ADR-037 choice is \textbf{bolded}.}
\label{tab:depth-compare}
\begin{tabular}{lllr}
\toprule
\(\varepsilon\) & \(r\) & Derivation & \(d\) \\
\midrule
\(2^{-23}\) (float32 machine eps) & \(0.99\) & Gate T software default & 1587 \\
\(10^{-6}\) & \(0.99\) & Earlier TableGen spec (erroneous) & 1380 \\
\(2^{-23}\) & \(0.9\) & float32 eps, hw bound & 152 \\
\(\mathbf{2^{-20}}\) & \(\mathbf{0.9}\) & \textbf{ADR-037 normative} & \textbf{132} \\
\(10^{-6}\) & \(0.9\) & Earlier hw\_dispatch (erroneous magic 150) & 132$^*$ \\
\(10^{-5}\) & \(0.9\) & Relaxed convergence & 110 \\
\bottomrule
\multicolumn{4}{l}{$^*$The magic number 150 in earlier code is inconsistent with both this row (132) and} \\
\multicolumn{4}{l}{\phantom{$^*$}the float32 eps row (152). It has no derivation.}
\end{tabular}
\end{table}

\subsection{Tightness: When \(\dhw = 132\) is Not Sufficient}

\begin{proposition}[Depth Insufficiency for Near-Boundary Sessions]
\label{prop:tightness}
If \(\rho(\Wp) = \rhw - \delta\) for small \(\delta > 0\), then the
absolute error after \(\dhw\) steps satisfies:
\[
  \abs{\rhohat_{\dhw} - \rho(\Wp)} \lesssim \rho(\Wp)^{\dhw}
  \approx (0.9 - \delta)^{132} \approx \epshw \cdot e^{132\delta/0.9}.
\]
For \(\delta = 0\) (boundary session, \(\rho(\Wp) = 0.9\)):
\[
  \abs{\rhohat_{132} - 0.9} \lesssim 0.9^{132} = e^{132\ln 0.9}
  \approx 2.78 \times 10^{-7} < \epshw.
\]
The pipeline is tight at the boundary: error is exactly \(\epshw\) to
leading order.
\end{proposition}

%% ═══════════════════════════════════════════════════════════════════════════
\section{Irreducibility Certificate: Algebraic Invariant and Complexity}
\label{sec:irred}
%% ═══════════════════════════════════════════════════════════════════════════

\subsection{Irreducibility via Rank}

\begin{definition}[Irreducible Weight Matrix]
\label{def:irred}
A weight matrix \(\Wp \in \R^{n \times n}\) is \emph{irreducible at eigenvalue
\(\lambda_1\)} if the dominant eigenspace is one-dimensional, i.e.,
\[
  \dim\ker(\Wp - \lambda_1 \Id) = 1.
\]
Equivalently, \(\rank(\Wp - \lambda_1 \Id) = n - 1\).
\end{definition}

\begin{theorem}[Irreducibility Certificate via Rank Check]
\label{thm:irred}
Let \(\hat\lambda_1 = \rhohat_{\dhw}\) be the hardware spectral estimate.
Under the gap condition \(\Delta\lambda \geq \dgap\) (so that
\cref{eq:correct} holds), define the perturbed matrix:
\[
  E = \Wp - \hat\lambda_1 \Id.
\]
Then:
\[
  \rank(E) = n - 1
  \;\Longleftrightarrow\;
  \dim\ker(\Wp - \lambda_1 \Id) = 1,
\]
with error probability zero over the continuum of real matrices (i.e., the
condition \(\rank(E) < n-1\) occurs on a measure-zero set of matrices for
any \(\hat\lambda_1 \neq \lambda_i\) for \(i \geq 2\)).
\end{theorem}

\begin{proofbox}
Since \(\hat\lambda_1 \approx \lambda_1\) (by \cref{thm:correct}),
\(E = \Wp - \hat\lambda_1 \Id \approx \Wp - \lambda_1 \Id\).
The kernel of \(\Wp - \lambda_1 \Id\) is the \(\lambda_1\)-eigenspace.
If \(\dim\ker(\Wp - \lambda_1 \Id) = 1\), then \(\rank(\Wp - \lambda_1 \Id) = n-1\).

For the estimated version, write \(\hat\lambda_1 = \lambda_1 + \epsilon\)
where \(\abs{\epsilon} \leq \epshw \cdot \rho(\Wp)\) by \cref{eq:correct}. Then:
\[
  E = (\Wp - \lambda_1\Id) - \epsilon \Id.
\]
The rank of \(E\) equals \(n-1\) if and only if exactly one singular value
of \((\Wp - \lambda_1\Id)\) has magnitude \(\leq \abs{\epsilon}\). Since
\(\Wp - \lambda_1\Id\) has one zero singular value (corresponding to
\(\mathbf{u}_1\)) and the next smallest is at least
\(\abs{\lambda_1 - \lambda_2} \geq \dgap \cdot \abs{\lambda_1} \geq \dgap \cdot \epshw\)
(much larger than \(\abs{\epsilon} \leq \epshw \cdot \rhw\)), the rank
of \(E\) equals \(n-1\) whenever \(\epshw \cdot \rhw < \dgap \cdot \epshw
\cdot \abs{\lambda_1}\), which reduces to \(\rhw < \dgap \cdot \abs{\lambda_1}\),
satisfied for all sessions with \(\abs{\lambda_1} > 0\) and
\(\dgap = 0.01\), \(\rhw = 0.9\). \hfill\(\square\)
\end{proofbox}

\subsection{Computational Complexity}

\begin{proposition}[Irreducibility Check Complexity]
\label{prop:irred-complexity}
The irreducibility check \(\rank(\Wp - \rhohat_{\dhw}\,\Id) \stackrel{?}{=} n-1\)
via truncated SVD has complexity:
\begin{equation}
  \mathcal{O}(n^2 \cdot \min(n, k_\varepsilon))
  \quad\text{where}\quad
  k_\varepsilon = \mathcal{O}\!\left(\log(1/\varepsilon)/\Delta\lambda\right).
  \label{eq:complexity}
\end{equation}
For the Tier 0 reference implementation using \texttt{numpy.linalg.matrix\_rank},
the full SVD complexity is \(\mathcal{O}(n^3)\). However, since only the
smallest singular value needs to be distinguished from zero, the randomized
one-shot SVD (Halko--Martinsson--Tropp, 2011) gives \(\mathcal{O}(n^2)\).
\end{proposition}

\begin{remark}
In the hardware pipeline, the irreducibility check is performed at Tier 0
(Python, before dispatch) and the result is cached in the digital twin
snapshot. It does not occur on the FPGA or ASIC datapath. The on-die
Prime Router implicitly assumes irreducibility for sessions that have passed
the certificate check.
\end{remark}

%% ═══════════════════════════════════════════════════════════════════════════
\section{Operator Norm Bounds for Composite Session Identifiers}
\label{sec:composite}
%% ═══════════════════════════════════════════════════════════════════════════

Composite session identifiers \(m = p_i \cdot p_j\) arise at inter-session
composition boundaries. The corresponding weight matrix is:

\begin{definition}[Composite Session Weight Matrix]
For coprime \(p_i, p_j \in \Pcal\), the composite weight matrix is:
\[
  W^{(p_i \cdot p_j)} = W^{(p_i)} \otimes W^{(p_j)} + \Phi^{(p_i, p_j)},
\]
where \(\otimes\) denotes the Kronecker product and \(\Phi^{(p_i,p_j)}\)
is the inter-session coupling term (ADR-SHF-003 morphism residual).
\end{definition}

\begin{theorem}[Composite Operator Norm Bound]
\label{thm:composite}
Under the decoupling assumption \(\norm{\Phi^{(p_i,p_j)}}_2 \leq \phi\) for
some coupling constant \(\phi \geq 0\):
\begin{equation}
  \rho(W^{(p_i \cdot p_j)})
  \leq \rho(W^{(p_i)}) \cdot \rho(W^{(p_j)}) + \phi.
  \label{eq:composite-rho}
\end{equation}
For the hardware contractivity invariant to hold at the composite level, it
is sufficient that:
\begin{equation}
  \rho(W^{(p_i)}) \cdot \rho(W^{(p_j)}) + \phi \leq \rhw.
  \label{eq:composite-hw}
\end{equation}
\end{theorem}

\begin{proofbox}
By the spectral radius bound for Kronecker products:
\[
  \rho(A \otimes B) = \rho(A) \cdot \rho(B).
\]
By subadditivity of the spectral radius under perturbation
(Weyl's inequality for non-Hermitian matrices, using the spectral norm
perturbation bound \(\abs{\rho(A + E) - \rho(A)} \leq \norm{E}_2\)):
\[
  \rho(W^{(p_i)} \otimes W^{(p_j)} + \Phi)
  \leq \rho(W^{(p_i)} \otimes W^{(p_j)}) + \norm{\Phi}_2
  = \rho(W^{(p_i)}) \cdot \rho(W^{(p_j)}) + \phi.
\]
\cref{eq:composite-hw} follows immediately from \cref{eq:composite-rho}
and the hardware invariant. \hfill\(\square\)
\end{proofbox}

\begin{corollary}[Independent Session Contractivity]
\label{cor:independent}
If \(\phi = 0\) (decoupled sessions) and both prime sessions are contractive
with \(\rho(W^{(p_i)}), \rho(W^{(p_j)}) \leq \rhw\), then:
\[
  \rho(W^{(p_i \cdot p_j)}) \leq \rhw^2 = 0.81 < \rhw.
\]
Decoupled composite sessions are automatically more contractive than their
components.
\end{corollary}

\begin{remark}
\cref{cor:independent} justifies the ASIC design choice that the inter-session
crossbar is needed only for non-zero \(\phi\) (coupled sessions). For purely
decoupled compositions, the prime-per-lane architecture is sufficient.
\end{remark}

%% ═══════════════════════════════════════════════════════════════════════════
\section{Contractivity Gap Lemma}
\label{sec:gap}
%% ═══════════════════════════════════════════════════════════════════════════

The contractivity gap lemma quantifies the safety margin between the true
spectral radius and the hardware bound, accounting for float32 round-off in
the spectral pipeline.

\begin{lemma}[Contractivity Gap under Round-Off]
\label{lem:gap}
Let \(\Wp\) be a weight matrix with \(\rho(\Wp) \leq \rhw - \Delta\)
for some gap \(\Delta > 0\). Let \(\tilde\rhohat_d\) be the spectral
estimate computed in float32 arithmetic (machine epsilon \(\mu = 2^{-24}\)).
Then the hardware contractivity verdict is correct (i.e.,
\(\tilde\rhohat_d < \rhw\)) provided:
\begin{equation}
  \Delta > \rhw \cdot \epshw + \rhw \cdot \dhw \cdot \mu \cdot \kappa(\Wp),
  \label{eq:gap}
\end{equation}
where \(\kappa(\Wp) = \norm{\Wp}_2 \cdot \norm{\Wp^{-1}}_2\) is the
condition number of \(\Wp\).
\end{lemma}

\begin{proofbox}
The total error in \(\tilde\rhohat_d\) has two sources:
\begin{enumerate}[noitemsep]
  \item \textit{Algorithmic error:} \(\abs{\rhohat_d - \rho(\Wp)} \leq \epshw \cdot \rhw\)
        by \cref{thm:correct} (assuming gap condition).
  \item \textit{Round-off error:} Each of the \(\dhw\) matrix-vector
        multiplications accumulates relative round-off of order \(n \mu\) in
        float32. Over \(\dhw\) steps, the accumulated error in the norm
        estimate is bounded by \(\rhw \cdot \dhw \cdot \mu \cdot \kappa(\Wp)\)
        (standard backward error analysis for power iteration in
        finite-precision arithmetic).
\end{enumerate}
The total error satisfies:
\[
  \abs{\tilde\rhohat_d - \rho(\Wp)}
  \leq \rhw \cdot \epshw + \rhw \cdot \dhw \cdot \mu \cdot \kappa(\Wp).
\]
For the verdict \(\tilde\rhohat_d < \rhw\) to be correct, we need
\(\rho(\Wp) + \abs{\tilde\rhohat_d - \rho(\Wp)} < \rhw\), i.e.,
\(\rhw - \rho(\Wp) > \abs{\tilde\rhohat_d - \rho(\Wp)}\), which gives
\cref{eq:gap}. \hfill\(\square\)
\end{proofbox}

\begin{hwbox}
\cref{lem:gap} implies that sessions with \(\rho(\Wp)\) very close to \(\rhw\)
(gap \(\Delta \approx 0\)) risk incorrect hardware verdicts due to round-off.
The hardware enforces \(\rhw = 0.9\) precisely to maintain a safety margin
below the software bound of \(r = 0.99\). For well-conditioned matrices
(\(\kappa(\Wp) \leq 100\)) with float32 \(\mu = 2^{-24} \approx 6 \times 10^{-8}\):
\[
  \Delta > 0.9 \times 2^{-20} + 0.9 \times 132 \times 6\times 10^{-8} \times 100
  \approx 8.6 \times 10^{-7} + 7.1 \times 10^{-4}
  \approx 7.1 \times 10^{-4}.
\]
Sessions must have spectral radius at least \(7.1 \times 10^{-4}\) below
\(\rhw = 0.9\) for guaranteed correct hardware verdicts.
\end{hwbox}

%% ═══════════════════════════════════════════════════════════════════════════
\section{Governance Commitment Hash: Collision Resistance}
\label{sec:hash}
%% ═══════════════════════════════════════════════════════════════════════════

\begin{definition}[Tier Commitment Hash]
\label{def:hash}
The commitment hash is:
\begin{equation}
  C = H\!\left(\mathrm{tier} \;\|\; t \;\|\; p_i \;\|\; \epshw \;\|\; \rhw \;\|\; \dhw\right),
  \label{eq:hash-full}
\end{equation}
where \(H\) is SHA-256 and \(\|\) denotes byte-level concatenation.
\end{definition}

\begin{proposition}[Constant-Drift Detection]
\label{prop:hash}
Let \(\Kcal = (\epshw, \rhw, \dhw)\) and \(\tilde\Kcal = (\tilde\epshw, \tilde\rhw, \tilde\dhw)\)
be two distinct constant sets (e.g., current ADR-037 vs.\ a stale version
with \(\tilde\dhw = 150\)). Then, under the collision resistance of SHA-256:
\begin{equation}
  \Pr\!\left[H(\mathrm{tier}\|\,t\|\,p_i\|\,\Kcal) = H(\mathrm{tier}\|\,t\|\,p_i\|\,\tilde\Kcal)\right]
  \leq 2^{-128},
  \label{eq:collision}
\end{equation}
i.e., the probability that a stale-constant bitstream presents a valid
commitment hash is negligible (\(\leq 2^{-128}\) under SHA-256's birthday
bound for 256-bit outputs).
\end{proposition}

\begin{proofbox}
Since \(\Kcal \neq \tilde\Kcal\), the input payloads
\(m = \mathrm{tier}\|\,t\|\,p_i\|\,\Kcal\) and
\(\tilde m = \mathrm{tier}\|\,t\|\,p_i\|\,\tilde\Kcal\) differ.
By the collision resistance of SHA-256 (no known attack below
\(2^{128}\) work), \(\Pr[H(m) = H(\tilde m)] \leq 2^{-128}\). \hfill\(\square\)
\end{proofbox}

\begin{remark}
The commitment hash \cref{eq:hash-full} serves a dual purpose: (1) it
authenticates the tier state transition (existing ADR-032 function), and
(2) it pins the hardware constant set at the time of bitstream compilation
(ADR-037 amendment). Any drift in \(\Kcal\)---even a single-bit change
to \(\dhw\)---produces a detectably different hash. This makes ADR-037
constant drift a cryptographically enforced invariant.
\end{remark}

%% ═══════════════════════════════════════════════════════════════════════════
\section{Summary of Norm Bounds and Hardware Implications}
\label{sec:summary}
%% ═══════════════════════════════════════════════════════════════════════════

\begin{table}[H]
\centering
\caption{Summary of operator norm bounds and their hardware implications}
\label{tab:summary}
\begin{tabular}{p{3.8cm} p{5cm} p{4.5cm}}
\toprule
\textbf{Result} & \textbf{Bound} & \textbf{Hardware Implication} \\
\midrule
Operator norm decay (\cref{thm:decay})
  & \(\norm{(\Wp)^k}_2 \leq C_p \cdot r^k\)
  & Drives pipeline depth formula; \(C_p = 1\) for normal matrices \\[4pt]
Uniform bound (\cref{prop:uniform})
  & \(\sup_p \norm{(\Wp)^k}_2 \leq C_{\mathcal{W}} r_0^k\)
  & Justifies fixed-depth pipeline across all sessions \\[4pt]
Power iteration error (\cref{thm:poweriter})
  & \(\abs{\rhohat_k - \abs{\lambda_1}} \leq \abs{\lambda_1} C_{\mathrm{iter}} \Delta\lambda^{-2} (|\lambda_2/\lambda_1|)^{2k}\)
  & Gap condition \(\Delta\lambda \geq \dgap\) required for correctness \\[4pt]
Pipeline correctness (\cref{thm:correct})
  & \(\abs{\rhohat_{\dhw} - \rho(\Wp)} \leq \epshw \cdot \rho(\Wp)\)
  & \texttt{CONTRACTIVE} verdict correct when \(\rho < \rhw(1-\epshw)\) \\[4pt]
Canonical depth (\cref{prop:depth})
  & \(\dhw = \lceil \log\epshw / \log\rhw \rceil = 132\)
  & Single source of truth; replaces magic number 150 \\[4pt]
Irreducibility (\cref{thm:irred})
  & \(\rank(\Wp - \rhohat\Id) = n-1\)
  & Required before ASIC block-diagonal routing \\[4pt]
Composite sessions (\cref{thm:composite})
  & \(\rho(W^{(p_ip_j)}) \leq \rho(W^{(p_i)})\rho(W^{(p_j)}) + \phi\)
  & Crossbar only needed for \(\phi > 0\) \\[4pt]
Round-off gap (\cref{lem:gap})
  & \(\Delta > \rhw\epshw + \rhw\,\dhw\,\mu\,\kappa(\Wp)\)
  & Safety margin \(\approx 7.1 \times 10^{-4}\) for \(\kappa \leq 100\) \\[4pt]
Hash collision (\cref{prop:hash})
  & \(\Pr[\text{collision}] \leq 2^{-128}\)
  & Stale constants undetectable only with \(< 2^{-128}\) probability \\
\bottomrule
\end{tabular}
\end{table}

%% ═══════════════════════════════════════════════════════════════════════════
\section*{Notation Index}
\addcontentsline{toc}{section}{Notation Index}

\begin{table}[H]
\centering
\begin{tabular}{ll}
\toprule
\textbf{Symbol} & \textbf{Meaning} \\
\midrule
\(\epshw\) & Hardware convergence threshold \(= 2^{-20}\) (ADR-037) \\
\(\rhw\) & Hardware contraction ceiling \(= 0.9\) (ADR-037) \\
\(\dhw\) & Canonical spectral pipeline depth \(= 132\) (ADR-037) \\
\(\dgap\) & Eigenvalue gap threshold \(= 0.01\) (ADR-037) \\
\(\Pmax\) & Largest 32-bit prime \(= 2^{31}-1\) (ADR-037) \\
\(\rho(A)\) & Spectral radius of \(A\) \\
\(\rhohat_d\) & Power iteration estimate of spectral radius after \(d\) steps \\
\(\norm{A}_2\) & Spectral (operator) 2-norm of \(A\) \\
\(\norm{A}_F\) & Frobenius norm of \(A\) \\
\(\kappa(A)\) & Condition number \(\norm{A}_2 \norm{A^{-1}}_2\) \\
\(\Delta\lambda\) & Relative eigenvalue gap \((\abs{\lambda_1}-\abs{\lambda_2})/\abs{\lambda_1}\) \\
\(\Vp\) & Prime-indexed multiplicity space \(\R^n\) \\
\(\Wp\) & Session weight matrix for prime \(p\) \\
\(\Kcal\) & Hardware constants triple \((\epshw,\rhw,\dhw)\) \\
\(C_p\) & Contractivity gap constant for session \(p\) \\
\(C_{\mathrm{iter}}\) & Power iteration initial alignment constant \\
\(\mu\) & Float32 machine epsilon \(= 2^{-24}\) \\
\(H\) & SHA-256 hash function \\
\(C\) & Tier commitment hash \\
\bottomrule
\end{tabular}
\end{table}


\nocite{*}
\bibliographystyle{unsrt}  
\bibliography{references}
\end{document}